\subsection{Harmonic synthesis}

Let A be a commutative Banach algebra. Recall that if M is a subset of \(\mathsf{X}(\mathsf{A})\) , we denote by \(\Upsilon(\mathsf{M})\) the intersection of the kernels of the elements of M (cf.~I, p.~30); this is an ideal of A.

PROPOSITION 4. --- Let A be a regular commutative Banach algebra without radical. Let F be a closed subset of X(A). The set of ideals I of A such that V(I) = F, ordered by inclusion, admits a greatest element, namely Y(F), and a least element, namely the set J of x ∈ A such that G(x) has compact support disjoint from F.

By construction, \(\Upsilon(F)\) is an ideal of A such that \(V(\Upsilon(F))\) contains F, and it is the largest ideal of A having this property. Since F is closed, there exists an ideal I of A such that \(V(I) = F\) ; we therefore have \(I \subset \Upsilon(F)\) , whence \(V(\Upsilon(F)) \subset V(I) = F\) so that \(V(\Upsilon(F)) = F\) . This proves the first assertion.

The set J is an ideal of A and V(J) contains F. Let us show that V(J) = F. Let \(\chi \in \mathsf{X}(\mathsf{A})\) not belong to F. Let U be a compact neighborhood of \(\chi\) not meeting F (TG, I, p.~65, cor. of prop. 9). By assertion (ii) of prop. 1 of I, p.~88, there exists \(x \in A\) such that \(\mathcal{G}(x)\) is equal to 1 at \(\chi\) and to 0 outside of U. We then have \(x \in J\) and therefore \(\chi \notin \mathrm{V}(\mathrm{J})\) . This shows that V(J) ⊂ F and hence V(J) = F.

Finally, let I be an ideal of A such that \(V(I) = F\) . We show that \(J \subset I\) . Let \(x \in J\) and let C be the support of \(\mathcal{G}(x)\) ; the set C is a compact subset of \(X(A)\) disjoint from F. By Prop. 3, there exists an element \(u \in I\) such that \(\mathcal{G}(u) = 1\) on C. We then have \(\mathcal{G}(x) = \mathcal{G}(ux)\) and therefore x = ux since A is radical-free (prop. 8 of I, p.~38). Consequently, we have \(x \in I\) , which shows that \(J \subset I\) .

COROLLARY 1. --- Let A be a regular commutative Banach algebra without radical. Let J be the set of \(x \in \mathrm{A}\) such that \(\mathcal{G}(x)\) has compact support. We assume that \(\overline{\mathrm{J}} = \mathrm{A}\) . Then every closed ideal of A distinct from A is contained in a regular maximal ideal.

If I is a closed ideal of A which is contained in no regular maximal ideal, then \(V(I) = \varnothing\) , hence \(I \supset J\) (prop. 4 applied to \(F = \varnothing\) ), whence \(I \supset \overline{J} = A\) .

COROLLARY 2. --- Let A be a commutative regular Banach algebra without radical. Let \(x, y \in A\) . If the support of \(\mathcal{G}(x)\) is compact and contained in the set of characters \(\chi\) such that \(\mathcal{G}(y)(\chi) \neq 0\) , then \(x\) is a multiple of \(y\) in \(A\) .

Let I be the ideal Ay of A. Then V(I) is the set of zeros of \(\mathcal{G}(y)\) . Since the support F of \(\mathcal{G}(x)\) is compact and disjoint from V(I), we have \(x \in I\) (prop. 4 applied to F).

DEFINITION 2. --- Let A be a commutative Banach algebra.

Let I be an ideal of A, \(x \in A\) , and \(\chi \in \mathsf{X}'(\mathsf{A})\) . We say that x belongs to I in a neighborhood of \(\chi\) if there exists an element \(y \in I\) such that \(\mathcal{G}'(y)\) and \(\mathcal{G}'(x)\) coincide in a neighborhood of \(\chi\) .

We say that A satisfies the Ditkin condition if, for every \(\chi \in \mathsf{X}'(\mathrm{A})\) and every \(x \in \mathrm{A}\) such that \(\mathcal{G}'(x)\) vanishes at \(\chi\) there exists a sequence \((x_n)\) in A such that \(x = \lim_{n \to \infty} x_n x\) and such that each \(\mathcal{G}'(x_n)\) vanishes in a neighborhood \(\mathrm{V}_n\) of \(\chi\) .

Remarks. --- Let A be a commutative Banach algebra.

\begin{enumerate}
\def\labelenumi{\arabic{enumi})}
\item
  If \(\chi\) is such that \(\mathcal{G}'(x)\) vanishes in a neighborhood of \(\chi\) , then \(x\) belongs to I in a neighborhood of \(\chi\) .
\item
  If x belongs to I in a neighborhood of \(\chi\) and \(y \in A\) is an arbitrary element, then xy belongs to I in a neighborhood of \(\chi\) .
\item
  The set of \(\chi\) such that \(x\) belongs to I in a neighborhood of \(\chi\) is open in \(\mathrm{X}'(\mathrm{A})\) .
\item
  Suppose that A is regular and radical-free. If x belongs to I in a neighborhood of \(\chi\) for all \(\chi \in \mathsf{X}'(\mathrm{A})\) , then x belongs to I (cor. 2 of I, p.~91 applied to the function \(f = \mathcal{G}'(x)\) and prop. 8 of I, p.~38).
\item
  Suppose that A is regular. Let I be an ideal of A and let \(\chi\) be an element of \(\mathsf{X}(\mathsf{A})\) such that \(\chi\notin\mathsf{V}(\mathsf{I})\) . Then every element \(x\) of A belongs to I in the neighborhood of \(\chi\) . Indeed, by def. 1 of I, p.~89, there exists a \(z\in\mathsf{A}\) such that \(\mathcal{G}'(z)\) is equal to 1 in a neighborhood of \(\chi\) and equal to 0 in a neighborhood of V(I). The support of \(\mathcal{G}(z)\) is compact and hence we have \(z\in\mathsf{I}\) (prop. 4 applied to V(I)); therefore \(xz\in\mathsf{I}\) , and \(\mathcal{G}'(xz)=\mathcal{G}'(x)\) in a neighborhood of \(\chi\) .
\end{enumerate}

Recall that a subspace K of a topological space X is said to be perfect if it is closed and has no isolated point (TG, I, p.~8).

Lemma 1. --- Let A be a commutative regular Banach algebra without radical, satisfying Ditkin's condition. Let I be a closed ideal of A and let x be an element of \(\Upsilon(\mathrm{V}(\mathrm{I}))\) . Let K be the set of \(\chi \in \mathsf{X}'(\mathrm{A})\) such that x does not belong to I in a neighborhood of \(\chi\) . Then the set K is a perfect subset of \(\mathsf{X}'(\mathrm{A})\) .

Let G denote the complement of K in \(X'(A)\) . The set G is open in \(X'(A)\) (remark 3), hence K is closed.

Proceed by contradiction, and suppose that K has an isolated point \(\chi_{0}\) . Let U be a neighborhood of \(\chi_{0}\) such that \(U - \{\chi_{0}\}\subset G\) . Since x does not belong to I in a neighborhood of \(\chi_{0}\) , Remark 5 shows that \(\chi_{0}\in\mathrm{V(I)}\) . In particular, we have \(\chi_{0}(x)=0\) since \(x\in\Upsilon(\mathrm{V(I)})\) .

We are going to show that there exists an element \(y\) of A which belongs to I in a neighborhood of every point of \(\mathsf{X}'(\mathrm{A}) - \{\chi_0\}\) , which does not belong to I in a neighborhood of \(\chi_0\) , and such that \(\chi_0(y) = 0\) .

Suppose first that the existence of such an element \(y\) has been proved. Since A satisfies Ditkin's condition, there then exists a sequence \((x_n)\) in A such that \(x_n y\) tends to \(y\) and such that each \(\mathcal{G}'(x_n)\) vanishes in a neighborhood of \(\chi_0\) . For every \(n\) , the element \(x_n y\) then belongs to I in a neighborhood of every point of \(X'(A)\) (Remarks 1 and 2) and therefore \(x_{n}y \in I\) (remark 4). Since I is closed, we deduce that \(y \in I\) , which contradicts the fact that y does not belong to I in a neighborhood of \(\chi_{0}\) .

It remains to prove the existence of \(y\) . If \(\chi_0 \neq 0\) , by assertion (iii) of prop. 1 of I, p.~88, there exists a \(u \in A\) such that \(\mathcal{G}'(u)\) is equal to 1 in a neighborhood of \(\chi_0\) and equal to 0 in a neighborhood of \(\mathrm{X}'(\mathrm{A}) - \mathrm{U}\) . Let \(y = ux\) . Since \(x\) belongs to I in a neighborhood of \(\chi\) for all \(\chi \in \mathrm{U} - \{\chi_0\}\) , the same is true of \(y\) . Moreover, if \(\chi \in \mathrm{X}'(\mathrm{A}) - \mathrm{U}\) , then \(\mathcal{G}'(y)\) vanishes in a neighborhood of \(\chi\) . Therefore (Remark 5) the element \(y = ux\) belongs to I in a neighborhood of every \(\chi \neq \chi_0\) . Since \(\mathcal{G}'(y)\) coincides with \(\mathcal{G}'(x)\) in a neighborhood of \(\chi_0\) , the fact that \(\chi_0\) belongs to K implies that \(y\) does not belong to I in a neighborhood of \(\chi_0\) . Finally, we have \(\chi_0(y) = \chi_0(u)\chi_0(x) = 0\) .

If \(\chi_0 = 0\) , there similarly exists an element \(v \in \mathrm{A}\) such that \(\mathcal{G}'(v)\) is zero in a neighborhood of \(\chi_0\) and equal to 1 in a neighborhood of \(\mathsf{X}'(\mathsf{A}) - \mathsf{U}\) ; as before, we deduce from this that the element \(y = x - vx\) belongs to I in a neighborhood of every \(\chi \neq \chi_0\) , that it does not belong to I in a neighborhood of \(\chi_0\) , and that \(\chi_0(y) = 0\) .

Lemma 2. — Let X be a topological space. Let F and D be disjoint subspaces of X such that F is closed and D is discrete. If F contains no nonempty perfect subspace, then neither does F ∪ D.

Indeed, suppose that K is a nonempty perfect subspace of F ∪ D. Let x be a point of K. If x belongs to D, it is isolated in D, hence also in F ∪ D since F is closed. Therefore x is isolated in K, which contradicts the hypotheses.

PROPOSITION 5. — Let A be a regular commutative Banach algebra without radical, satisfying Ditkin's condition. Let I be a closed ideal of A such that the boundary F of V(I) contains no nonempty perfect set. Then I = Y(V(I)), that is, I is the set of x ∈ A such that G(x) vanishes on V(I). In particular, if V(I) reduces to a single point χ, we have I = Ker(χ).

We have \(\mathrm{I} \subset \Upsilon(\mathrm{V}(\mathrm{I}))\) . Now let \(x \in \Upsilon(\mathrm{V}(\mathrm{I}))\) . Let \(\mathrm{G}\) be the set of characters \(\chi \in \mathsf{X}'(\mathbf{A})\) such that \(x\) belongs to \(\mathrm{I}\) in a neighborhood of \(\chi\) . It is open and its complement \(\mathrm{K}\) is perfect (lemma 1). Since \(\mathcal{G}'(x)\) is zero on \(\mathrm{V}(\mathrm{I})\) , the set \(\mathrm{G}\) contains the interior of \(\mathrm{V}(\mathrm{I}) \cup \{0\}\) (remark 1). It also contains \(\mathsf{X}(\mathbf{A}) - \mathsf{V}(\mathbf{I})\) by remark 5. Therefore \(\mathrm{K} = \mathsf{X}'(\mathbf{A}) - \mathsf{G}\) is contained in the boundary \(\mathrm{F}_0\) of \(\mathrm{V}(\mathrm{I}) \cup \{0\}\) .

We have \(F_{0} \subset F \cup \{0\}\) . The hypothesis therefore implies that \(F_{0}\) contains no nonempty perfect set (lemma 2). It follows therefore from lemma 1 that the perfect set K is empty. Hence x belongs to I in a neighborhood of every \(\chi \in \mathsf{X}'(\mathsf{A})\) , which means that \(x \in I\) (remark 4). Thus \(\Upsilon(\mathrm{V}(\mathrm{I})) \subset \mathrm{I}\) , which concludes the proof.

\section{STELLAR ALGEBRAS}

In this paragraph, the base field is C.

\subsection{Semilinear involutions}

Let E be a complex vector space. A semilinear involution on E is an R-linear map from E to E such that \(u \circ u = Id_{E}\) and \(u(\lambda x) = \overline{\lambda}u(x)\) for all \(\lambda \in C\) and every \(x \in E\) . We then denote by \(E^{u}\) the real vector subspace of E formed by the elements \(x \in E\) such that \(u(x) = x\) .

Lemma 1. --- Let E be a complex vector space and let u be a semilinear involution on E. Let \(x \in \mathrm{E}\) ; set

\[
x _ {1} = \frac {1}{2} (x + u (x)), \quad x _ {2} = \frac {1}{2 i} (x - u (x)).
\]

The pair \((x_{1}, x_{2})\) is the unique element of \(E^{u} \times E^{u}\) such that \(x = x_{1} + ix_{2}\) . The elements \(x_{1}\) and \(x_{2}\) satisfy \(x_{1} + ix_{2} = x\) and belong to \(E^{u}\) since \(u(u(x)) = x\) . Conversely, if \(y_{1}\) and \(y_{2}\) are in \(E^{u}\) and satisfy \(x = y_{1} + iy_{2}\) , it follows that \(u(x) = u(y_{1}) + u(iy_{2}) = y_{1} - iy_{2}\) , hence

\[
y _ {1} = \frac {1}{2} (x + u (x)) = x _ {1}, \quad i y _ {2} = \frac {1}{2} (x - u (x)) = i x _ {2}.
\]

PROPOSITION 1. --- Let \(E_{1}\) and \(E_{2}\) be complex vector spaces, and let \(u_{1}\) and \(u_{2}\) be semilinear involutions on \(E_{1}\) and \(E_{2}\) respectively. Let f be a linear map from \(E_{1}\) into \(E_{2}\) . Then \(f \circ u_{1} = u_{2} \circ f\) if and only if \(f(\mathrm{E}_{1}^{u_{1}}) \subset \mathrm{E}_{2}^{u_{2}}\) .

If \(f \circ u_1 = u_2 \circ f\) , we immediately obtain \(f(\mathrm{E}_1^{u_1}) \subset \mathrm{E}_2^{u_2}\) . Conversely, suppose this condition is satisfied. Let \(x \in \mathrm{E}\) . Write \(x = x_1 + ix_2\) with \((x_1, x_2) \in \mathrm{E}_1^{u_1} \times \mathrm{E}_1^{u_1}\) (lemma above). We then have \(f(u_1(x)) = f(x_1) - if(x_2)\) and \(u_2(f(x)) = u_2(f(x_1)) - iu_2(f(x_2)) = f(x_1) - if(x_2) = f(u_1(x))\) .

\subsection{Involutive algebras}

DEFINITION 1. --- Let A be an algebra over C. An involution on A is a map \(x \mapsto x^{*}\) from A to A such that:

\[
\begin{array}{c} (x ^ {*}) ^ {*} = x, \quad (x + y) ^ {*} = x ^ {*} + y ^ {*}, \quad (\lambda x) ^ {*} = \overline {{\lambda}} x ^ {*} \\ (x y) ^ {*} = y ^ {*} x ^ {*} \end{array}
\]

for all \(x, y \in A\) and \(\lambda \in C\) . An algebra over C equipped with an involution is called an involutive algebra.

An involution on A is in particular an isomorphism of the ring A onto the opposite ring A°.

Let A be an involutive algebra. One says that \(x^{*}\) is the adjoint of x. A subset of A stable under the involution is said to be self-adjoint. If A has a unit element e, one has \(e^{*} = e\) ; one says that (A, e) is a unital involutive algebra. An element u of a unital involutive algebra is said to be unitary if \(uu^{*} = u^{*}u = e\) , in other words if u is invertible and its inverse is \(u^{*}\) .

An element \(x \in A\) is said to be hermitian if \(x = x^{*}\) and normal if \(xx^{*} = x^{*}x\) . This terminology generalizes that of A, IX, § 7, no. 3. Every hermitian element is normal, and every unitary element is normal. The set \(A_{h}\) of hermitian elements of A is a real vector subspace of A. If x and y are hermitian and commute, one has \((xy)^{*} = y^{*}x^{*} = yx = xy\) , hence xy is hermitian. For every \(x \in A\) , the elements \(xx^{*}\) and \(x^{*}x\) of A are hermitian.

If A = C equipped with the involution \(z \mapsto \overline{z}\) , one has \(A_{h} = R\) .

Lemma 2. --- Let A be an involutive algebra and \(x \in \mathrm{A}\) . The elements

\[
x _ {1} = \frac {1}{2} (x + x ^ {*}), \qquad x _ {2} = \frac {1}{2 i} (x - x ^ {*})
\]

are hermitian and satisfy \(x = x_{1} + ix_{2}\) . If \(x = y_{1} + iy_{2}\), where \(y_{1}\) and \(y_{2}\) are hermitian, then \(x_{1} = y_{1}\) and \(x_{2} = y_{2}\) . Moreover, the element x is normal if and only if \(x_{1}\) and \(x_{2}\) commute.

The first two assertions follow from lemma 1 of I, p.~95. One computes that \(xx^{*} - x^{*}x = 2i(x_{2}x_{1} - x_{1}x_{2})\) , hence x is normal if and only if \(x_{1}\) and \(x_{2}\) commute.

Let A be a unital involutive algebra. For \(x \in A\) to be invertible, it is necessary and sufficient that \(x^{*}\) be so, and one then has \((x^{*})^{-1} = (x^{-1})^{*}\) . Since \((x - \lambda e)^{*} = x^{*} - \overline{\lambda}e\) for all \(\lambda \in C\) , one deduces that \(\mathrm{Sp}_{\mathrm{A}}(x^{*}) = \overline{\mathrm{Sp}_{\mathrm{A}}(x)}\) .

Let A be an involutive algebra and \(\widetilde{\mathrm{A}}\) the algebra obtained from A by adjoining a unit element. There exists on \(\widetilde{\mathrm{A}}\) a unique involution extending that of A, given by \((\lambda, x)^{*} = (\overline{\lambda}, x^{*})\) for \(\lambda \in \mathbf{C}\) and \(x \in \mathrm{A}\) . If \(x \in \mathrm{A}\) , one has \(\mathrm{Sp}_{\mathrm{A}}'(x^{*}) = \overline{\mathrm{Sp}_{\mathrm{A}}'(x)}\) .

Let A and B be involutive algebras. A morphism from A to B is an algebra morphism \(\varphi\) from A to B such that \(\varphi(x^{*}) = \varphi(x)^{*}\) for all \(x\) in A. The identity map of A is a morphism of involutive algebras. If C is an involutive algebra and \(\pi: B \to C\) a morphism of involutive algebras, then \(\pi \circ \varphi\) is a morphism of involutive algebras. If \(\varphi\) is an isomorphism of algebras, then \(\varphi^{-1}\) is a morphism of involutive algebras, and one says that \(\varphi\) is an isomorphism of involutive algebras.

By prop. 1 of I, p.~95, if A and B are involutive algebras, an algebra morphism \(\varphi\) from A to B is a morphism of involutive algebras if and only if \(\varphi(A_h) \subset B_h\) . One calls an involutive subalgebra of A a self-adjoint subalgebra. The center of A is an involutive subalgebra. If \(A_1\) is a self-adjoint two-sided ideal of A, the involution of A defines by passage to the quotient an involution on the algebra \(A/A_1\) , and the canonical map from A onto \(A/A_1\) is a morphism of involutive algebras.

Let A be an involutive algebra. The radical of A is equal to the radical of the opposite algebra (A, VIII, p.~431, prop. 7), and is therefore self-adjoint.

Let A be an involutive algebra. If M ⊂ A is self-adjoint, its commutant M' is an involutive subalgebra of A. If \(x \in A\) , the bicommutant of \(\{x, x^{*}\}\) is an involutive subalgebra containing x and \(x^{*}\) and this subalgebra is commutative if and only if x is normal (no. 5 of I, p. 5).

Remark. — Let A be an involutive algebra and let B be a maximal commutative involutive subalgebra of A. Then B is a maximal commutative subalgebra. In particular, if A is unital, then the algebra B is full.

Indeed, let \(x \in A\) be an element commuting with B. Then \(x^{*}\) commutes with B. Write \(x = x_{1} + ix_{2}\) with \(x_{1}\) and \(x_{2}\) hermitian; the elements \(x_{1}\) and \(x_{2}\) commute with B (lemma 2). The subalgebra of A generated by B and \(x_{1}\) is therefore commutative and involutive. Consequently, it is equal to B, so that \(x_{1} \in B\) . Similarly, we have \(x_{2} \in B\) , and hence finally \(x \in B\) .

Let A be an involutive algebra. If \(f\) is a linear form on A, the map \(x \mapsto \overline{f(x^*)}\) on A is a linear form on A, which is denoted by \(f^*\) . The map \(f \mapsto f^*\) is a semilinear involution on A'. We say that \(f\) is Hermitian if \(f = f^*\). By lemma 1 of I, p.~95, every linear form \(f\) on A has a unique representation \(f = f_1 + if_2\) where \(f_1\) and \(f_2\) are Hermitian, namely \(f_1 = \frac{1}{2}(f + f^*)\) and \(f_2 = \frac{1}{2i}(f - f^*)\) .

For a linear form \(f\) to be Hermitian, it is necessary and sufficient that the restriction of \(f\) to \(A_h\) be real-valued (proposition 1 of I, p.~95). The map \(f \mapsto f|A_h\) is an isomorphism of the real vector space of Hermitian forms on the dual vector space of the real vector space \(A_h\) .

In particular, we denote by \(\mathsf{X}^{\prime}(\mathsf{A})_{h}\) (resp. \(\mathsf{X}(\mathsf{A})_{h}\) ) the set of Hermitian characters of A (resp. the set of nonzero Hermitian characters of A). A character is therefore Hermitian if its restriction to \(A_{h}\) takes real values.

If A is commutative and if \(\chi\) is a character of A, then \(\chi^{*}\) is a character of A, and the map \(\chi \mapsto \chi^{*}\) is a homeomorphism from \(\mathsf{X}'(\mathsf{A})\) onto \(\mathsf{X}'(\mathsf{A})\) .

Lemma 3. --- Let A be a commutative involutive algebra. The Gelfand transform of A to \(\mathcal{C}_{0}(\mathsf{X}(\mathsf{A}))\) is a morphism of involutive algebras if and only if every character of A is Hermitian.

Indeed, saying that \(G_{A}\) is a morphism of involutive algebras amounts to saying that, for all \(x \in A\) and \(\chi \in \mathsf{X}(A)\) , one has

\[
\chi (x ^ {*}) = \mathcal {G} _ {\mathrm{A}} (x ^ {*}) (\chi) = \overline {{\mathcal {G} _ {\mathrm{A}} (x) (\chi)}} = \overline {{\chi (x)}},
\]

that is, every \(\chi\) is Hermitian.

Examples. --- 1) Let A be the algebra of complex-valued functions on a set X. The map \(f \mapsto \overline{f}\) is an involution in A. The subalgebra of bounded functions in X is an involutive subalgebra of A. If X is a locally compact topological space, the subalgebras \(\mathcal{C}(X), \mathcal{C}_b(X), \mathcal{C}_0(X)\) and \(\mathcal{K}(X)\) are involutive subalgebras of A.

\begin{enumerate}
\def\labelenumi{\arabic{enumi})}
\setcounter{enumi}{1}
\item
  Let X be a locally compact topological space and \(\mu\) a positive measure on X. The map \(f \mapsto \overline{f}\) is an involution on the algebra \(\mathcal{L}^{\infty}(X, \mu)\) ; it induces, by passing to the quotient, an involution on the unital algebra \(\mathrm{L}^{\infty}(X, \mu)\) .
\item
  Let E be a complex Hilbert space. On the Banach algebra \(\mathcal{L}(\mathrm{E})\) , the map \(x \mapsto x^{*}\) (EVT, V, p.~37, prop. 1) is an involution.
\item
  Let G be a locally compact group. Let \(\mathcal{M}^{1}(G)\) be the Banach algebra of bounded complex measures on G (example 6 of I, p.~19).
\end{enumerate}

The map \(x \mapsto x^{-1}\) from G to G maps every measure \(\mu \in \mathcal{M}^{1}(G)\) into a measure \(\check{\mu} \in \mathcal{M}^{1}(G)\) (INT, VII, p.~12, formula (13)). We denote by \(\mu^{*}\) the complex conjugate measure of \(\check{\mu}\) . The map \(\mu \mapsto \check{\mu}\) is an isometric isomorphism from the Banach algebra \(\mathcal{M}^{1}(G)\) onto the Banach algebra \(\mathcal{M}^{1}(G^{\circ})\) (INT, VIII, §3, no. 1, cor. of prop. 7); hence \(\mu \mapsto \mu^{*}\) is an isometric involution of the Banach algebra \(\mathcal{M}^{1}(G)\) .

The set A of bounded measures admitting a density with respect to a Haar measure is a closed subalgebra of \(\mathcal{M}^{1}(\mathrm{G})\) stable under the involution (cf. INT, VIII, §4, no. 5); it does not depend on the choice of a Haar measure.

Let \(\nu\) be a left Haar measure on G and denote by \(\Delta\) the module of G. We equip \(\mathrm{L}^1 (\mathrm{G},\nu)\) with the product \((f,g)\mapsto f*\nu g\) and with the involution \(f\mapsto f^{*} = \widetilde{f}\cdot \Delta^{-1}\) , where \(\widetilde{f} (x) = \overline{f(x^{-1})}\) for all \(x\in \mathbf{G}\). Then the map \(f\mapsto f\cdot \nu\) is an isomorphism of the involutive algebra \(\mathrm{L}^1 (\mathrm{G},\nu)\) onto A. This isomorphism is isometric. In particular, \(\mathrm{L}^1 (\mathrm{G},\nu)\) is identified with an involutive subalgebra of \(\mathcal{M}^1 (\mathrm{G})\) .

\begin{enumerate}
\def\labelenumi{\arabic{enumi})}
\setcounter{enumi}{4}
\tightlist
\item
  Let U be an open subset of C stable under complex conjugation. Consider the algebra \(\mathcal{O}(\mathrm{U})\) of complex-valued holomorphic functions on U. For every function \(f \in \mathcal{O}(\mathrm{U})\) , the map \(f^{*}: z \mapsto \overline{f(\overline{z})}\) is a holomorphic function on U. The map \(f \mapsto f^{*}\) is an involution on \(\mathcal{O}(\mathrm{U})\) .
\end{enumerate}

Similarly, let S be a compact subset of C stable under complex conjugation. Consider the algebra \(\mathcal{O}(S)\) of germs of complex-valued holomorphic functions in a neighborhood of S. The map \(f \mapsto f^{*}\) is an involution on \(\mathcal{O}(S)\) .

The subalgebra \(\mathcal{O}_{\mathbf{R}}(\mathrm{U})\) (resp. the subalgebra \(\mathcal{O}_{\mathbf{R}}(\mathrm{S})\) ) defined in no. 13 of I, p.~85 is the set of Hermitian elements of \(\mathcal{O}(\mathrm{U})\) (resp. of \(\mathcal{O}(\mathrm{S})\) ).

\subsection{Normed involutive algebras}

DEFINITION 2. --- A normed involutive algebra is a normed algebra A equipped with an involution \(x \mapsto x^{*}\) such that \(\|x^{*}\| = \|x\|\) for all \(x\) . If A is a Banach algebra, then A is said to be an involutive Banach algebra.

Examples. --- 1) Let X be a locally compact topological space. The Banach algebra \(\mathcal{C}_{b}(\mathrm{X})\) of continuous and bounded complex functions on X, equipped with the norm \(\|f\| = \sup_{x\in\mathrm{X}}|f(x)|\) and with the involution \(f\mapsto\overline{f}\) , is an involutive Banach algebra. The subalgebra \(\mathcal{C}_{0}(\mathrm{X})\) of continuous functions tending to 0 at infinity is a closed involutive subalgebra of \(\mathcal{C}_{b}(\mathrm{X})\) .

\begin{enumerate}
\def\labelenumi{\arabic{enumi})}
\setcounter{enumi}{1}
\item
  Let X be a locally compact topological space and \(\mu\) a positive measure on X. The involutive algebra \(\mathrm{L}^{\infty}(\mathrm{X},\mu)\) (example 2 of I, p.~99) is an involutive Banach algebra, since \(|f| = |\overline{f}|\) for every element \(f \in \mathrm{L}^{\infty}(\mathrm{X},\mu)\) .
\item
  The involutive algebra \(\mathcal{L}(\mathrm{E})\) of continuous endomorphisms of a complex Hilbert space E (I, p.~99, example 3), equipped with the norm
\end{enumerate}

\[
\| u\| = \sup_{\substack{x\in \operatorname{E}\\ \| x\| \leqslant 1}}\| u(x)\|
\]

(EVT, III, p.~14) is an involutive Banach algebra (EVT, V, p.~37, prop. 1).

\begin{enumerate}
\def\labelenumi{\arabic{enumi})}
\setcounter{enumi}{3}
\item
  The involutive algebra \(\mathcal{M}^{1}(G)\) of bounded measures on a locally compact group (I, p.~99, example 4), equipped with the usual norm (example 6 of I, p.~19), is an involutive Banach algebra. Let \(\nu\) be a left Haar measure on G. The involutive Banach algebra \(\mathrm{L}^{1}(\mathrm{G},\nu)\) identifies with a closed subalgebra of \(\mathcal{M}^{1}(\mathrm{G})\) .
\item
  Let \((\mathrm{A}_{i})\) be a family of involutive normed algebras. Let A be the product normed algebra of the \(A_{i}\) (No. 1 of I, p.~15). The algebra A, equipped with the involution \((x_{i})^{*} = (x_{i}^{*})\) , is an involutive normed algebra. If each of the algebras \(A_{i}\) is an involutive Banach algebra, then A is an involutive Banach algebra. We say that A is the product involutive normed algebra (resp. the product involutive Banach algebra) of \(A_{i}\) .
\item
  Let A be a normed involutive algebra, and let \(\widetilde{\mathrm{A}}\) be the normed algebra deduced from A by adjoining a unit element. Equipped with the involution defined in no. 2, the algebra \(\widetilde{\mathrm{A}}\) is a normed involutive algebra. If A is an involutive Banach algebra, then \(\widetilde{\mathrm{A}}\) is also an involutive Banach algebra.
\end{enumerate}

If A is a normed involutive algebra, the closure of an involutive subalgebra is an involutive subalgebra. If M ⊂ A, the smallest closed involutive subalgebra containing M is called the closed involutive subalgebra generated by M; it is the closure of the subalgebra generated by M ∪ M*. If M consists of a single normal element, the closed involutive subalgebra generated by M is commutative, and all its elements are normal.

Similarly, if A is a unital involutive normed algebra and M a subset of A, the smallest closed unital involutive subalgebra containing M is called the closed unital involutive subalgebra generated by M; it is the closure of the unital subalgebra generated by M ∪ M*. If M consists of a single normal element, the closed unital involutive subalgebra generated by M is commutative, and all its elements are normal.

The quotient of an involutive normed algebra by a closed two-sided self-adjoint ideal, the completion, and the opposite of an involutive normed algebra are naturally involutive normed algebras.

If A is a normed involutive algebra, the set \(A_{h}\) of the Hermitian elements of A is a real normed vector space.

Lemma 4. --- Let A be a normed involutive algebra. For every continuous linear form f on A, we have \(\|f^{*}\|=\|f\|\) . If moreover f is Hermitian, then \(\|f\|=\|f|A_{h}\|\) .

The first assertion follows from the definitions. For the second, let \(g\) be the restriction of \(f\) to \(\mathrm{A}_h\) . We have \(\| f \| \geqslant \| g \|\) . Let us prove the reverse inequality. For every \(\varepsilon > 0\) , there exists \(x \in \mathrm{A}\) such that \(\| x \| \leqslant 1\) and \(|f(x)| \geqslant \|f\| - \varepsilon\) . By multiplying x by a complex number of modulus 1, we may assume \(f(x) \geqslant 0\) . Then the element \(\frac{1}{2}(x + x^{*})\) belongs to \(A_{h}\) and has norm \(\leqslant 1\) . We have

\[
\left| g \left(\frac {1}{2} (x + x ^ {*})\right) \right| = \frac {1}{2} | f (x) + f (x ^ {*}) | = f (x) \geqslant \| f \| - \varepsilon
\]

hence \(\|g\|\geqslant\|f\|-\varepsilon\) . From this we deduce that \(\|g\|\geqslant\|f\|\) .

In what follows, we shall identify continuous Hermitian linear forms on A with continuous real linear forms on \(A_{h}\) .

Lemma 5. --- Let A be an involutive Banach algebra.

\begin{enumerate}
\def\labelenumi{\alph{enumi})}
\item
  For every \(x \in A\) , one has \(\exp(x)^{*} = \exp(x^{*})\) ;
\item
  Let \(x \in \mathrm{A}_h\) be a Hermitian element. Then \(\exp(ix)\) is unitary.
\end{enumerate}

Indeed, since the involution on A is continuous, we have

\[
\exp (x) ^ {*} = \left(\sum_ {n = 0} ^ {\infty} \frac {x ^ {n}}{n !}\right) ^ {*} = \sum_ {n = 0} ^ {\infty} \frac {(x ^ {*}) ^ {n}}{n !} = \exp (x ^ {*})
\]

for all \(x \in A\) (formula (9) of I, p.~78). If \(x \in A_{h}\) , it then follows that

\[
\exp (i x) ^ {*} = \exp ((i x) ^ {*}) = \exp (- i x) = \exp (i x) ^ {- 1}
\]

(formula (11) of I, p.~78).

\subsection{Stellar algebras}

DEFINITION 3. --- A stellar algebra is an involutive Banach algebra A such that \(\|x\|^{2}=\|x^{*}x\|\) for all \(x\in A\) .

If A and B are stellar algebras, a morphism, or morphism of stellar algebras, from A to B is a morphism of involutive algebras from A to B. An isomorphism from A to B is an isomorphism of involutive algebras from A to B.

Some authors speak of "C*-algebra".

Let A be a stellar algebra. A stellar subalgebra of A is a closed involutive subalgebra of A.

Examples. --- 1) The involutive Banach algebra of continuous endomorphisms of a complex Hilbert space (example 3 of I, p.~100) is a stellar algebra (EVT, V, p.~39, prop. 2).

\begin{enumerate}
\def\labelenumi{\arabic{enumi})}
\setcounter{enumi}{1}
\tightlist
\item
  Let X be a locally compact topological space. The involutive Banach algebra \(\mathcal{C}_{b}(X)\) of continuous bounded complex-valued functions on X (example 1 of I, p.~100) is a stellar algebra. Indeed, for every function \(f \in \mathcal{C}_b(\mathrm{X})\) , one has \(f^*f = |f|^2\) , and hence \(\| f^*f \| = |||f|^2 \| = \| f\|^2\) .
\end{enumerate}

Let \(\mathrm{A} = \mathcal{C}_0(\mathrm{X})\) be the involutive Banach subalgebra of functions tending to 0 at infinity. It is a stellar subalgebra of \(\mathcal{C}_b(\mathrm{X})\) . For every function \(f \in \mathrm{A}\) , one has \(\| f \| = \varrho(f)\) since \(\mathrm{Sp}_{\mathrm{A}}'(f) = f(\mathrm{X}) \cup \{0\}\) .

Let X and Y be locally compact topological spaces. For every proper partial map \(\varphi\) from X into Y (def. 1 of I, p.~33), the algebra morphism \(\varphi^{*}\) of \(\mathcal{C}_{0}(\mathrm{Y})\) into \(\mathcal{C}_{0}(\mathrm{X})\) (prop. 3 of I, p.~34) is a morphism of involutive algebras. Conversely, every morphism of stellar algebras \(\pi\colon\mathcal{C}_{0}(\mathrm{Y})\to\mathcal{C}_{0}(\mathrm{X})\) is of this form (loc. cit.).

\begin{enumerate}
\def\labelenumi{\arabic{enumi})}
\setcounter{enumi}{2}
\item
  Let X be a compact topological space and \(x_{0} \in X\) a fixed element of X. The subalgebra \(\mathcal{C}'(X)\) of \(\mathcal{C}(X)\) of continuous functions \(f: X \to C\) such that \(f(x_{0}) = 0\) is a stellar algebra.
\item
  Let X be a separated topological space and \(\mu\) a positive measure on X. The involutive Banach algebra \(\mathrm{L}^{\infty}(\mathrm{X},\mu)\) is a commutative unital stellar algebra.
\item
  Let \((\mathbf{A}_{i})\) be a family of stellar algebras. The product involutive Banach algebra A of the \(A_{i}\) (example 5 of I, p.~101) is a stellar algebra, called the product stellar algebra of the \(A_{i}\) .
\item
  Let A be a stellar algebra. If B is a closed involutive subalgebra of A, then B is a stellar algebra. It will be seen (V, to appear) that every stellar algebra is isomorphic to a closed involutive subalgebra of the stellar algebra of endomorphisms of a Hilbert space (example 1).
\item
  Let A be a stellar algebra. If M ⊂ A is any subset, then the closed involutive subalgebra of A generated by M is a stellar algebra, called the stellar subalgebra of A generated by M. If A is moreover unital, then the closed unital involutive subalgebra generated by M is a unital stellar algebra, called the unital stellar subalgebra of A generated by M.
\item
  In general, the involutive Banach algebra \(\mathcal{M}^{1}(G)\) (example 4 of I, p.~100) is not a stellar algebra (I, p.~181, exerc. 8).
\end{enumerate}

Lemma 6. --- Let A be a Banach algebra equipped with an involution satisfying

\[
\| x \| ^ {2} \leqslant \| x ^ {*} x \|\tag{1}
\]

for all \(x \in \mathbf{A}\) . Then A is a stellar algebra.

Let \(x \in \mathbf{A}\) . We then have \(\|x\|^2 \leqslant \|x^*\| \cdot \|x\|\) , whence \(\|x\| \leqslant \|x^*\|\) . By exchanging the roles of \(x\) and \(x^*\) we see that \(\|x\| = \|x^*\|\) . Thus \(\|x^*x\| \leqslant \|x^*\|\|x\| \leqslant \|x\|^2\) and the hypothesis implies the equality \(\|x\|^2 = \|x^*x\|\) .

Lemma 7. --- Let A be a stellar algebra.

\begin{enumerate}
\def\labelenumi{\alph{enumi})}
\tightlist
\item
  The regular representation \(\gamma\) of A (def. 1 of I, p.~16) is isometric, that is,
\end{enumerate}

\[
\| x \| = \sup _ {\| y \| \leqslant 1} \| x y \|,
\]

for all \(x\in \mathbf{A}\) .

\begin{enumerate}
\def\labelenumi{\alph{enumi})}
\setcounter{enumi}{1}
\tightlist
\item
  For every \(x \in A_{h}\) , one has
\end{enumerate}

\[
\varrho (x) = \| x \|.\tag{2}
\]

Let \(x \in \mathrm{A}\) . We have \(\sup_{\| y \| \leqslant 1} \| xy \| \leqslant \| x \|\) . To prove that \(\| x \| \leqslant \sup_{\| y \| \leqslant 1} \| xy \|\) , we may assume \(\| x \| = 1\) . Then the element \(y = x^*\) satisfies \(\| y \| = \| x^* \| = 1\) , and \(\| xy \| = \| x \|^2 = 1\) , whence the assertion \(a\) ).

Suppose that \(x\) is Hermitian. It follows \(\| x^2 \| = \| x^* x \| = \| x \|^2\) , whence by induction \(\| x^{2^n} \|^{2^{-n}} = \| x \|\) for every integer \(n \geqslant 1\) , whence the assertion \(b)\) by prop. 1 of I, p.~20.

Remarks. --- 1) Let A be a unital stellar algebra. Then we have

\[
\left\| 1 \right\| ^ {2} = \left\| 1 ^ {*} 1 \right\| = \left\| 1 \right\|,
\]

therefore the norm \(\|1\|\) is zero or equal to 1. If A ≠ \{0\}, we deduce \(\|1\| = 1\) . Consequently, for every unitary element u of A, we have \(\|u\| = \|u^{*}u\|^{1/2} = 1\) .

\begin{enumerate}
\def\labelenumi{\arabic{enumi})}
\setcounter{enumi}{1}
\tightlist
\item
  Let A be a normed involutive algebra. If \(\|x\|^{2}=\|x^{*}x\|\) for all \(x\in A\) , the completion \(\widehat{A}\) of A is a stellar algebra.
\end{enumerate}

PROPOSITION 2. --- Let A be an involutive Banach algebra, let B be a stellar algebra, and let π be a morphism of involutive algebras from A to B. Then one has \(\|\pi(x)\| \leqslant \|x\|\) for all \(x \in A\) , and in particular π is continuous.

For every \(x \in \mathrm{A}\) , one has \(\mathrm{Sp}_{\mathrm{B}}^{\prime}(\pi(x)) \subset \mathrm{Sp}_{\mathrm{A}}^{\prime}(x)\) , hence \(\varrho(\pi(x)) \leqslant \varrho(x) \leqslant \|x\|\) . Since \(\pi(x^{*}x) \in \mathrm{B}_{h}\) , one has \(\|\pi(x^{*}x)\| = \varrho(\pi(x^{*}x))\) (formula (2)), hence

\[
\| \pi (x) \| ^ {2} = \| \pi (x ^ {*} x) \| = \varrho (\pi (x ^ {*} x)) \leqslant \| x ^ {*} x \| = \| x \| ^ {2}.
\]

PROPOSITION 3. --- Let A be a stellar algebra and let \(\widetilde{\mathrm{A}}\) be the involutive algebra obtained from A by adjoining a unit element. There exists a unique norm on \(\widetilde{\mathrm{A}}\) extending that of A and making \(\widetilde{\mathrm{A}}\) a stellar algebra.

The uniqueness of such a norm follows from Prop. 2. We now show its existence. Let us denote by \(\widetilde{e}\) the identity element of \(\widetilde{\mathrm{A}}\) .

First, suppose that A has a unit element e. The product of the involutive normed algebras A and \(\mathbf{C}(\widetilde{e}-e)\) is identified with \(\widetilde{A}\) and is a stellar algebra (example 5). The norm on \(\widetilde{A}\) extends that of A, whence the assertion.

Now suppose that A has no identity element. For every \(x \in \widetilde{A}\) , let \(\gamma_{x}\) be the multiplication operator \(y \mapsto xy\) from A to A, and set \(\|x\|_{\widetilde{A}} = \| \gamma_{x}\|\) . The map \(x \mapsto \|x\|_{\widetilde{A}}\) is a semi-norm on \(\widetilde{A}\) . For all x and \(x'\) in \(\widetilde{A}\) , one has \(\|xx'\|_{\widetilde{A}} \leqslant \|x\|_{\widetilde{A}} \|x'\|_{\widetilde{A}}\) . Moreover, by Lemma 7, we have \(\|x\|_{\widetilde{A}} = \|x\|\) for all \(x \in \widetilde{A}\) .

Let us show that the map \(x \mapsto \|x\|_{\widetilde{\mathrm{A}}}\) is a norm on \(\widetilde{\mathrm{A}}\) . Let \(\lambda \in \mathbf{C}\) and \(x \in \mathrm{A}\) such that \(\| \lambda \widetilde{e} + x\|_{\widetilde{\mathrm{A}}} = 0\) . If \(\lambda \neq 0\) the condition \((\lambda \widetilde{e} + x)y = 0\) for all \(y \in \mathrm{A}\) implies that \(-\lambda^{-1}x\) is a left identity element in A. Likewise, the element \(-\overline{\lambda}^{-1}x^{*}\) is a right identity element. Thus, the algebra A would then have an identity element, contrary to the hypothesis. Hence, \(\lambda = 0\) . But then \(0 = \|x\|_{\widetilde{\mathrm{A}}} = \|x\|\) , and hence \(x = 0\) .

Since A is complete and of codimension 1 in \(\widetilde{\mathrm{A}}\) , the space \(\widetilde{\mathrm{A}}\) equipped with the norm \(x \mapsto \|x\|_{\widetilde{\mathrm{A}}}\) is complete. To conclude, it is therefore sufficient to show that one has \(\|x\|_{\widetilde{\mathrm{A}}}^{2} \leqslant \|x^{*}x\|_{\widetilde{\mathrm{A}}}\) for all \(x \in \widetilde{\mathrm{A}}\) (lemma 6). One may assume that \(\|x\|_{\widetilde{\mathrm{A}}} = 1\) . For every real number \(r < 1\) , there therefore exists \(y \in \mathrm{A}\) such that \(\|y\| = \|y^{*}\| \leqslant 1\) and \(\|xy\|^2 \geqslant r\) . Since \(xy \in \mathrm{A}\) , one has

\[
\| x ^ {*} x \| _ {\widetilde {A}} \geqslant \| x ^ {*} x y \| \geqslant \| y ^ {*} (x ^ {*} x) y \| = \| (x y) ^ {*} (x y) \| = \| x y \| ^ {2} \geqslant r.
\]

From this one deduces that \(\| x^{*}x\|_{\widetilde{\mathrm{A}}}\geqslant 1\) , and hence the involutive algebra \(\widetilde{\mathrm{A}}\) endowed with the norm \(x\mapsto \| x\|_{\widetilde{\mathrm{A}}}\) is a stellar algebra.

DEFINITION 4. --- One says that \(\widetilde{\mathrm{A}}\) , endowed with the norm of prop. 3, is the stellar algebra deduced from A by adjoining an identity element.

When \(A \neq \{0\}\) , the stellar-algebra norm on the normed involutive algebra \(\widetilde{A}\) is not the one considered in example 6 of I, p.~101 (cf.~exercise 10 of I, p.~181).

PROPOSITION 4. --- Let A be a stellar algebra.

\begin{enumerate}
\def\labelenumi{\alph{enumi})}
\item
  If A has a unit element and u is a unitary element of A, then \(\mathrm{Sp}(u) \subset \mathbf{U}\) ;
\item
  If h is a Hermitian element of A, then \(\mathrm{Sp}'(h) \subset \mathbf{R}\) .
\end{enumerate}

We may assume that A is nonzero. Let us prove the assertion \(a\) ). Let \(u\) be a unitary element of A. We have \(\|u\|=\|u^{-1}\|=1\) (remark 1), hence \(\mathrm{Sp}(u)\subset\mathbf{U}\) (I, p.~26, cor. 3). To prove \(b\) ), we may assume that A is unital (prop. 3). Let \(h\) be a Hermitian element of A. Then \(\exp(ih)\) is unitary (Lemma 5 of I, p.~102). Thus, by Cor. 2 of I, p.~67 and \(a\) ), we have \(\exp(i\mathrm{Sp}(h))=\mathrm{Sp}(\exp(ih))\subset\mathbf{U}\) , which means that \(\mathrm{Sp}(h)\subset\mathbf{R}\) .

PROPOSITION 5. — Let A be a unital stellar algebra and B a stellar subalgebra of A containing the identity element of A. Then B is a full subalgebra of A. In particular, one has \(\mathrm{Sp}_{\mathrm{B}}(x) = \mathrm{Sp}_{\mathrm{A}}(x)\) for all \(x\) in B.

Let x be a Hermitian element of B. Since \(\mathrm{Sp}_{\mathrm{B}}(x) \subset \mathbf{R}\) (prop. 4), prop. 6 of I, p.~28 shows that \(\mathrm{Sp}_{\mathrm{B}}(x) = \mathrm{Sp}_{\mathrm{A}}(x)\) . In particular, x is invertible in B if and only if it is invertible in A.

Now let x be any element of B that is invertible in A. Then \(x^{*}\) is invertible in A, and \(xx^{*}\) is invertible in A. Since \(xx^{*}\) is Hermitian, the preceding shows that \(xx^{*}\) is invertible in B. This implies that x is right-invertible in B. Similarly, one verifies that x is left-invertible in B, and consequently that x is invertible in B. Thus, B is a full subalgebra of A.

COROLLARY. --- Let A be a stellar algebra and let B be a stellar subalgebra of A. Then we have \(\mathrm{Sp}_{\mathrm{B}}^{\prime}(x)=\mathrm{Sp}_{\mathrm{A}}^{\prime}(x)\) for all x in B.

By adjoining an identity element (prop. 3), this follows from prop. 5.

PROPOSITION 6 (Fuglede–Putnam Theorem)

Let A be a unital stellar algebra. Let a and b be normal elements of A. If \(c \in A\) satisfies ac = cb, then \(a^{*}c = cb^{*}\) .

The hypothesis implies \((wa)^{k}c = c(wb)^{k}\) for every integer \(k \geqslant 0\) and every \(w \in C\) , hence \(e^{wa}c = ce^{wb}\) (formula (9) of I, p.~78). Consider the function f from C to A defined by \(z \mapsto e^{-za^{*}}ce^{zb^{*}}\) . This is a holomorphic function on C, whose derivative satisfies

\[
f ^ {\prime} (z) = - a ^ {*} e ^ {- z a ^ {*}} c e ^ {z b ^ {*}} + e ^ {- z a ^ {*}} c b ^ {*} e ^ {z b ^ {*}}
\]

for all \(z \in \mathbf{C}\) . Since \(c = e^{-\overline{z}a} ce^{\overline{z}b}\) one can write

\[
f (z) = e ^ {- z a ^ {*}} e ^ {\overline {{z}} a} c e ^ {- \overline {{z}} b} e ^ {z b ^ {*}}.
\]

Since \(a\) and \(b\) are normal, the elements \(e^{-za^*}e^{\overline{z} a} = e^{-za^* + \overline{z} a}\) and \(e^{-\overline{z} b}e^{zb^*} = e^{-\overline{z} b + zb^*}\) of A are unitary for all \(z\in \mathbf{C}\) (Lemma 5 of I, p.~102), hence of norm 1. Consequently, one has \(\| f(z)\| \leqslant \| c\|\) for all \(z\in \mathbf{C}\) . The function \(f\) is therefore constant (VAR, R1, 3.3.6, p.~29), that is, \(f(z) = f(0) = c\) for all \(z\in \mathbf{C}\) . But then \(-a^{*}c + cb^{*} = f'(0) = 0\) .

COROLLARY. --- Let A be a stellar algebra and let a be a normal element of A. The commutant (resp. bicommutant) of \(\{a, a^{*}\}\) coincides with the commutant (resp. bicommutant) of a.

It suffices to take b = a in the proposition.

\subsection{Commutative stellar algebras}

Lemma 8. --- Let X and Y be metric spaces, with X complete. Let f be a mapping from X into Y such that

\[
d (f (x), f (y)) \geqslant d (x, y)
\]

for all x and y in X and such that the graph of f is closed in X × Y. Then f is a closed map.

Let F be a closed subset of X and let \((y_{n})_{n\in\mathbf{N}}\) a sequence in \(f(\mathrm{F})\) which converges to \(y\in Y\) . For every \(n\in N\) , let \(x_{n}\in F\) such that \(f(x_{n})=y_{n}\) . The hypothesis implies that the sequence \((x_{n})_{n\in\mathbf{N}}\) is a Cauchy sequence; let \(x\in X\) be its limit; it is an element of F, because F is closed. Moreover, we have \((x_{n},f(x_{n}))\to(x,y)\) into \(X\times Y\) . Since the graph of f is closed, it follows that \(y=f(x)\) belongs to \(f(\mathrm{F})\) .

Lemma 9. --- Let A be a commutative stellar algebra. Every character of A is Hermitian, and the Gelfand transformation is a morphism of stellar algebras from A into \(\mathcal{C}_{0}(\mathsf{X}(\mathsf{A}))\) .

It suffices to prove the first assertion (proposition 7 of I, p.~38 and lemma 3 of I, p.~98). Let \(\chi\) be a character of A. For every Hermitian element \(y\) of A, we have \(\chi(y) = \mathcal{G}_{\mathrm{A}}(y)(\chi) \in \mathrm{Sp}'(y) \subset \mathbf{R}\) (prop. 4 of I, p.~106). Consequently, the character \(\chi\) is Hermitian (prop. 1 of I, p.~95).

THEOREM 1. --- Let A be a commutative stellar algebra. The Gelfand transformation is an isometric isomorphism of the stellar algebra A onto the stellar algebra \(\mathcal{C}_{0}(\mathsf{X}(\mathsf{A}))\) of continuous functions on \(\mathsf{X}(\mathsf{A})\) vanishing at infinity.

The Gelfand transform is an involutive algebra morphism from A into \(\mathcal{C}_{0}(\mathsf{X}(\mathsf{A}))\) (lemma 9). Let B be its image. It is an involutive subalgebra of \(\mathcal{C}_{0}(\mathsf{X}(\mathsf{A}))\) . The elements of B separate the points of \(\mathsf{X}(\mathsf{A})\) by definition. Let \(\chi \in \mathsf{X}(\mathsf{A})\) . There exists \(x \in \mathsf{A}\) such that \(\chi(x) \neq 0\) , hence \(f = \mathcal{G}(x)\) is an element of B such that \(f(\chi) \neq 0\) . By TG, X, p.~40, cor. 2, the subalgebra B is therefore dense in \(\mathcal{C}_{0}(\mathsf{X}(\mathsf{A}))\) .

For every Hermitian element \(y\) of A, we have \(\| \mathcal{G}(y)\| = \varrho (y) = \| y\|\) (prop. 7 of I, p.~38 and formula (2) of I, p.~104), whence, for every \(x\in \mathrm{A}\) , the equalities

\[
\| x \| ^ {2} = \| x ^ {*} x \| = \| \mathcal {G} (x ^ {*} x) \| = \| \overline {{\mathcal {G} (x)}} \cdot \mathcal {G} (x) \| = \| \mathcal {G} (x) \| ^ {2}.
\]

Thus G is isometric, and its image B is therefore closed (Lemma 8). We conclude that \(\mathrm{B} = \mathcal{C}_{0}(\mathrm{X}(\mathrm{A}))\) .

COROLLARY 1. --- Let A be a stellar algebra and let x be a normal element of A. Then \(\|x\| = \varrho(x)\) .

Since the stellar subalgebra of A generated by x and \(x^{*}\) is commutative, we may assume that A is commutative. In this case, the corollary follows from Th. 1, Example 2 of I, p.~102, and Th. 1 of I, p.~24.

COROLLARY 2. --- Let A be a commutative stellar algebra.

\begin{enumerate}
\def\labelenumi{\alph{enumi})}
\item
  There exists a locally compact topological space X such that A is isomorphic to the stellar algebra \(\mathcal{C}_{0}(\mathrm{X})\) ;
\item
  Let B be a commutative stellar algebra. The map \(\pi \mapsto \mathsf{X}'(\pi)\) is a bijection from the set of stellar algebra morphisms from A to B onto the set of proper partial maps from \(\mathsf{X}(\mathsf{B})\) into \(\mathsf{X}(\mathsf{A})\) (def. 1 of I, p.~33).
\end{enumerate}

Theorem 1 establishes the first assertion, and the second assertion follows from Prop. 3 of I, p.~34 and from Example 2 of I, p.~102.

COROLLARY 3. --- Let A be a unital commutative stellar algebra.

\begin{enumerate}
\def\labelenumi{\alph{enumi})}
\item
  There exists a compact topological space X such that A is isomorphic to the stellar algebra \(\mathcal{C}(\mathrm{X})\) ;
\item
  Let B be a commutative unital stellar algebra. The map \(\pi \mapsto \mathsf{X}(\pi)\) is a bijection from the set of unital morphisms of stellar algebras from A to B to the set of continuous maps from \(\mathsf{X}(\mathsf{B})\) into \(\mathsf{X}(\mathsf{A})\) .
\end{enumerate}

This follows from the foregoing and from prop. 4 of I, p.~35.

Remark. --- *Let G be the category whose objects are locally compact spaces and whose morphisms are proper partial maps (def. 1 of I, p.~33), and let S be the category of commutative stellar algebras, whose morphisms are morphisms of involutive algebras. Consider the functor from S to the opposite category G° which associates to a commutative stellar algebra A the locally compact space X(A) of nonzero characters of A, and to a morphism φ: A → B of commutative stellar algebras the continuous map X'(φ). Th. 1 and Cor. 2 mean that this functor is an equivalence of categories, and that a quasi-inverse of this functor is the functor which associates to a locally compact topological space X the commutative stellar algebra C₀(X).

Similarly, Corollary 3 means that the opposite category of the category of compact spaces is equivalent to the category of unital commutative stellar algebras.*

\subsection{Functional calculus in unital stellar algebras}

In this number, A is a unital stellar algebra and x is a normal element of A.

Let B be the unital stellar subalgebra of A generated by x; it is commutative and contained in the bicommutant of \(\{x, x^{*}\}\) , hence in the bicommutant of x (corollary of Proposition 6 in I, p.~106). The Gelfand transform \(\mathcal{G}_{\mathrm{B}}: \mathrm{B} \to \mathcal{C}(\mathrm{X}(\mathrm{B}))\) is an isomorphism of stellar algebras (th. 1 of I, p.~108). We have \(\mathrm{Sp}_{\mathrm{B}}(x) = \mathrm{Sp}_{\mathrm{A}}(x)\) (prop. 5 of I, p.~106).

Lemma 10. --- The map \(\mathrm{ev}_x\colon \chi \mapsto \chi (x)\) induces a homeomorphism of \(\mathsf{X}(\mathbf{B})\) onto \(\mathrm{Sp}_{\mathrm{A}}(x)\) .

The map \(x \mapsto \chi(x)\) of \(\mathsf{X}(\mathsf{B})\) into \(\mathbf{C}\) is continuous, and its image is equal to \(\mathrm{Sp}_{\mathsf{B}}(x)\) by prop. 6 of I, p.~37, hence to \(\mathrm{Sp}_{\mathsf{A}}(x)\) . Since the characters of B are Hermitian (lemma 9 of I, p.~107), characters of B that agree on \(x\) are also equal on \(x^{*}\) and therefore are equal on the unital stellar subalgebra B generated by \(x\) . This proves that the map \(\mathrm{ev}_x\) is injective. Since \(\mathsf{X}(\mathsf{B})\) is compact and \(\mathbf{C}\) is separated, it induces a homeomorphism from \(\mathsf{X}(\mathsf{B})\) onto its image, hence the lemma.

From the lemma we deduce an isomorphism of stellar algebras \(\varphi_x\colon \mathcal{C}(\mathrm{Sp}_{\mathrm{A}}(x))\to \mathcal{C}(\mathrm{X(B)})\) . It maps a function \(f\in \mathcal{C}(\mathrm{Sp}_{\mathrm{A}}(x))\) to the function \(\chi \mapsto f(\chi (x))\) . The map \(\mathcal{G}_{\mathrm{B}}^{-1}\circ \varphi_{x}\) is an isomorphism of the stellar algebra \(\mathcal{C}(\mathrm{Sp}_{\mathrm{A}}(x))\) onto B.

DEFINITION 5. --- The involutive morphism \(f \mapsto (\mathcal{G}_{\mathrm{B}}^{-1} \circ \varphi_x)(f)\) of \(\mathcal{C}(\mathrm{Sp}_{\mathrm{A}}(x))\) to A is called the continuous functional calculus map of \(x\) in A. We denote it \(f \mapsto f(x)\) .

Remark. --- The map \(f \mapsto f(x)\) is isometric; its image is the unital stellar subalgebra B generated by x, which is contained in the bicommutant of x.

If \(f\) is the restriction to \(\mathrm{Sp}_{\mathrm{A}}(x)\) of a function of the form \(z \mapsto \mathrm{P}(z, \overline{z})\) , where \(\mathrm{P} \in \mathbf{C}[\mathrm{X}, \mathrm{Y}]\) is a polynomial, then we have \(f(x) = \mathrm{P}(x, x^{*})\) in the usual algebraic sense.

Examples. --- 1) Suppose that there exists \(\lambda \in \mathbf{C}\) such that \(\mathrm{Sp}_{\mathrm{A}}(x)\) is reduced to \(\lambda\) . We then have \(x = \lambda \cdot 1\) . Indeed, the identity function of \(\mathrm{Sp}_{\mathrm{A}}(x)\) is equal to \(\lambda\) , hence its image under the functional calculus map, that is, \(x\) is equal to \(\lambda \cdot 1\) .

\begin{enumerate}
\def\labelenumi{\arabic{enumi})}
\setcounter{enumi}{1}
\item
  For \(x\) to be Hermitian, it is necessary and sufficient that \(\mathrm{Sp}_{\mathrm{A}}(x)\) be contained in \(\mathbf{R}\) . Indeed, let \(f\) be the continuous map on \(\mathrm{Sp}_{\mathrm{A}}(x)\) given by \(f(z) = z - \overline{z}\) . Then \(x\) is Hermitian if and only if \(f(x) = 0\) , that is, if \(f\) is zero, that is, if \(\mathrm{Sp}_{\mathrm{A}}(x)\) is contained in \(\mathbf{R}\) .
\item
  For \(x\) to be unitary, it is necessary and sufficient that its spectrum be contained in the unit circle of \(\mathbf{C}\) . Indeed, let \(f \in \mathcal{C}(\mathrm{Sp}_{\mathrm{A}}(x))\) be the function defined by \(f(z) = z\overline{z} - 1\) ; the element \(x\) is unitary if and only if \(f(x) = 0\) , that is, if \(f\) is zero.
\end{enumerate}

Proposition 7. — The map \(f \mapsto f(x)\) is the unique unital morphism of involutive algebras from \(\mathcal{C}(\mathrm{Sp}_{\mathrm{A}}(x))\) into A such that the identity map \(z\) of \(\mathrm{Sp}_{\mathrm{A}}(x)\) has image \(x\) .

Indeed, the unital subalgebra of \(\mathcal{C}(\mathrm{Sp}_{\mathrm{A}}(x))\) generated by the elements \(z\) and \(\overline{z}\) of \(\mathcal{C}(\mathrm{Sp}_{\mathrm{A}}(x))\) is dense in \(\mathcal{C}(\mathrm{Sp}_{\mathrm{A}}(x))\) (TG, X, p.~40, cor. 1). Since every morphism of involutive algebras from \(\mathcal{C}(\mathrm{Sp}_{\mathrm{A}}(x))\) to A is continuous (I, p.~104, prop. 2), there exists at most one morphism of involutive algebras from \(\mathcal{C}(\mathrm{Sp}_{\mathrm{A}}(x))\) to A which maps \(z\) to \(x\) .

The following corollary shows that when \(f\) is the restriction of a holomorphic function in a neighborhood of \(\mathrm{Sp}_{\mathrm{A}}(x)\) the definition of \(f(x)\) coincides with that of the holomorphic functional calculus in one variable from no. 9 of I, p.~74.

COROLLARY 1. --- Let \(f \in \mathcal{O}(\mathrm{Sp}_{\mathrm{A}}(x))\) be a germ of a holomorphic function in a neighborhood of \(\mathrm{Sp}_{\mathrm{A}}(x)\) and let \(\widetilde{f} \in \mathcal{C}(\mathrm{Sp}_{\mathrm{A}}(x))\) be the continuous function on \(\mathrm{Sp}_{\mathrm{A}}(x)\) associated with \(f\) . We have \(\widetilde{f}(x) = f(x)\) , where \(f(x)\) is the element of A given by the holomorphic functional calculus.

Indeed, the map \(f \mapsto \widetilde{f}(x)\) is a unital continuous morphism from \(\mathcal{O}(\mathrm{Sp}_{\mathrm{A}}(x))\) to A which maps the germ of the identity function in a neighborhood of \(\mathrm{Sp}_{\mathrm{A}}(x)\) to \(x\) . The result then follows from th. 5 of I, p.~74.

COROLLARY 2. --- Let \(f \in \mathcal{C}(\mathrm{Sp}_{\mathrm{A}}(x))\) .

\begin{enumerate}
\def\labelenumi{\alph{enumi})}
\tightlist
\item
  We have
\end{enumerate}

\[
\operatorname{Sp} _ {\mathrm{A}} (f (x)) = f (\operatorname{Sp} _ {\mathrm{A}} (x));
\]

\begin{enumerate}
\def\labelenumi{\alph{enumi})}
\setcounter{enumi}{1}
\tightlist
\item
  For every \(g \in \mathcal{C}(\mathrm{Sp}_{\mathrm{A}}(f(x)))\) , one has \((g \circ f)(x) = g(f(x))\) .
\end{enumerate}

Since \(f(x)\) belongs to the full subalgebra B of A, we have \(\mathrm{Sp}_{\mathrm{A}}(f(x)) = \mathrm{Sp}_{\mathrm{B}}(f(x))\) (prop. 5 of I, p.~106). The isomorphism \(f \mapsto f(x)\) of \(\mathcal{C}(\mathrm{Sp}_{\mathrm{A}}(x))\) to B preserves the spectrum; we therefore have \(\mathrm{Sp}_{\mathrm{B}}(f(x)) = \mathrm{Sp}_{\mathcal{C}(\mathrm{Sp}_{\mathrm{A}}(x))}(f) = f(\mathrm{Sp}_{\mathrm{A}}(x))\) (example 3 of I, p.~17). This proves assertion \(a\) ).

The map \(g \mapsto (g \circ f)(x)\) is a unital morphism of involutive algebras from \(\mathcal{C}(\mathrm{Sp}_{\mathrm{A}}(f(x)))\) to A which maps the identity function of \(\mathrm{Sp}_{\mathrm{A}}(f(x))\) to \(f(x)\) . By prop. 7, we therefore have \((g \circ f)(x) = g(f(x))\) for all \(g \in \mathcal{C}(\mathrm{Sp}_{\mathrm{A}}(f(x)))\) .

Example 4. --- Let X be a locally compact space and let A = \(\mathcal{C}_{b}(X)\) be the commutative unital stellar algebra of continuous bounded functions on X (example 2 of I, p.~102). Let \(g \in A\) ; its spectrum S is the closure in C of the set \(g(\mathrm{X})\) of the values of g (example 3 of I, p.~17). The functional calculus mapping of g is then the mapping \(f \mapsto f \circ g\) , for \(f \in \mathcal{C}(\mathrm{S})\) . Indeed, this map is a unital morphism of stellar algebras such that the identity map of S has image g.

In the case where X is compact, we have \(A = \mathcal{C}(X)\) and \(S = g(X)\) .

Let \(\pi\colon\mathrm{A}\to\mathrm{A}^{\prime}\) be a unital morphism of unital stellar algebras. The element \(\pi(x)\) of \(\mathrm{A}^{\prime}\) is normal and its spectrum relative to \(\mathrm{A}^{\prime}\) is contained in \(\mathrm{Sp}_{\mathrm{A}}(x)\) . We then have:

PROPOSITION 8. --- Let \(f \in \mathcal{C}(\mathrm{Sp}_{\mathrm{A}}(x))\) . Still denoting \(f\) the restriction of \(f\) to \(\mathrm{Sp}_{\mathrm{A}'}(\pi(x))\) , we have the equality \(\pi(f(x)) = f(\pi(x))\) . In particular, for every \(\chi \in \mathsf{X}(\mathrm{A})\) , one has \(\chi(f(x)) = f(\chi(x))\) .

Let \(z\) be the identity map of \(\mathrm{Sp}_{\mathrm{A}}(x)\) . The maps defined by \(f \mapsto \pi(f(x))\) and \(f \mapsto f(\pi(x))\) are continuous unital morphisms of involutive algebras from \(\mathcal{C}(\mathrm{Sp}_{\mathrm{A}}(x))\) to B that map \(z\) to \(\pi(x)\) . These morphisms therefore coincide on the unital involutive subalgebra of \(\mathcal{C}(\mathrm{Sp}_{\mathrm{A}}(x))\) generated by \(z\) . Since this one is dense in \(\mathcal{C}(\mathrm{Sp}_{\mathrm{A}}(x))\) (TG, X, p.~40, cor. 1), these morphisms are equal.

COROLLARY. --- Suppose that A is commutative. For every \(f \in \mathcal{C}(\mathrm{Sp}_{\mathrm{A}}(x))\) , one has \(\mathcal{G}_{\mathrm{A}}(f(x)) = f \circ \mathcal{G}_{\mathrm{A}}(x)\) .

It suffices to apply Prop. 8 to the Gelfand transform \(\mathcal{G}_{\mathrm{A}}\) from A to \(\mathcal{C}(\mathsf{X}(\mathrm{A}))\) and to note (example above) that \(f(\mathcal{G}_{\mathrm{A}}(x)) = f \circ \mathcal{G}_{\mathrm{A}}(x)\) .

\subsection{Applications of functional calculus}

PROPOSITION 9. --- Every injective morphism of stellar algebras is isometric and, in particular, has closed image.

Let A and A′ be stellar algebras, and let \(\pi\colon A\to A'\) be a morphism of involutive algebras from A to A'.

Suppose first that A and A' are unital and that \(\pi\) is unital. We have \(\|\pi\| \leqslant 1\) (prop. 2). Suppose there exists x in A such that \(\|\pi(x)\| < \|x\|\) . Let \(y = x^{*}x\) ; it is a Hermitian element of A. Since A and A' are stellar algebras, we have \(\|\pi(y)\| = \|\pi(x)\|^{2} < \|x\|^{2} = \|y\|\) , that is, \(\varrho(\pi(y)) < \varrho(y)\) (lemma 7 of I, p.~104). In particular, \(\mathrm{Sp}_{\mathrm{A}'}(\pi(y))\) is a closed subset of \(\mathrm{Sp}_{\mathrm{A}}(y)\) , distinct from \(\mathrm{Sp}_{\mathrm{A}}(y)\)

(Remark 6 of I, p.~3 and Th. 1 of I, p.~24). There then exists a nonzero function \(f \in \mathcal{C}(\mathrm{Sp}_{\mathrm{A}}(y))\) such that \(f| \mathrm{Sp}_{\mathrm{A}'}(\pi(y)) = 0\) (TG, IX, p.~13, Prop. 3). Let \(w = f(y) \in \mathrm{A}\) . We have \(w \neq 0\) since \(f \neq 0\) but \(\pi(w) = \pi(f(y)) = f(\pi(y)) = 0\) since \(f\) is zero on \(\mathrm{Sp}_{\mathrm{A}'}(\pi(y))\) (Prop. 8). Hence \(\pi\) is not injective.

We now treat the general case. Let \(\widetilde{\mathrm{A}}\) and \(\widetilde{\mathrm{A}}'\) be the stellar algebras deduced from A and from \(\mathrm{A}'\) respectively by adjoining a unit element (def. 4 of I, p.~106). There exists a unique unital morphism of involutive algebras \(\widetilde{\pi}\colon \widetilde{\mathrm{A}}\to \widetilde{\mathrm{A}}'\) extending \(\pi\) . This morphism is injective, hence is isometric by what precedes. For every \(x\in \mathrm{A}\) , we then have \(\| \pi (x)\| = \| \widetilde{\pi} (x)\| = \| x\|\) .

Lemma 11. --- Let X be a completely regular topological space, that is, uniformizable and separated (TG, IX, p.~8, def. 4), containing at least two points. There exist nonzero continuous functions f and g in \(\mathcal{C}(\mathrm{X})\) such that \(fg = 0\) .

Let \(x \neq y\) be distinct points of X. Let U and V be open neighborhoods of x and y, respectively, such that U and V are disjoint. Since X is uniformizable, by TG, IX, p.~7, th. 2, there exists a function \(f \in \mathcal{C}(X)\) such that \(f(x) = 1\) and \(f|X - U = 0\) . Similarly, there exists \(g \in \mathcal{C}(X)\) such that \(g(y) = 1\) and \(g|X - V = 0\) . We then have fg = 0.

PROPOSITION 10. --- Let A be a unital stellar algebra. Suppose that for every pair \((x,y)\) of permutable elements of A, the condition xy = 0 implies that x = 0 or y = 0. Then A = C · 1.

If A is not equal to C·1, there exists a hermitian element x in A which does not belong to C·1 (lemma 2 of I, p.~96). Let B be the unital stellar subalgebra of A generated by x. It is commutative, and isomorphic to \(\mathcal{C}(\mathrm{Sp}_{\mathrm{A}}(x))\) (I, p.~110, remark). Since x is not scalar, its spectrum in B is not reduced to a single element (example 1 of I, p.~110). There therefore exist nonzero continuous functions f and g on \(\mathrm{Sp}_{\mathrm{A}}(x)\) such that fg = 0 (lemma 11). The elements \(f(x)\) and \(g(x)\) of A are nonzero, permutable, and satisfy \(f(x)g(x) = 0\) in A.

PROPOSITION 11. --- Let A be a unital stellar algebra, and let a, x and y be elements of A. Suppose that x and y are normal. If xa = ay, then we have \(f(x)a = af(y)\) for every continuous function f on the union of the spectrum of x and the spectrum of y. In particular, we have \(f(aa^{*})a = af(a^{*}a)\) for every function \(f \in \mathcal{C}(\mathrm{Sp}'(a^{*}a))\) .

Let \(\mathrm{S} = \mathrm{Sp}(x) \cup \mathrm{Sp}(y)\) . Proposition 6 of I, p.~106 implies that \(x^{*}a = ay^{*}\) . Consequently, we obtain \(f(x)a = af(y)\) for every function \(f\) which is of the form \(z \mapsto \mathrm{P}(z,\overline{z})\) , where \(\mathrm{P} \in \mathbf{C}[\mathrm{X},\mathrm{Y}]\) is a polynomial. Since the set of functions \(f \in \mathcal{C}(\mathrm{S})\) satisfying \(f(x)a = af(y)\) is a closed subalgebra of \(\mathcal{C}(\mathrm{S})\) , it coincides with \(\mathcal{C}(\mathrm{S})\) by TG, X, p.~40, cor. 1.

The second assertion is a consequence of the first, applied to the hermitian elements \(x = aa^{*}\) and \(y = a^{*}a\) , taking into account the fact that \(\mathrm{Sp}'(a^{*}a) = \mathrm{Sp}'(aa^{*})\) (prop. 1 of I, p.~5).

\subsection{Functional calculus in a non-unital algebra}

Let A be a stellar algebra and let \(\widetilde{\mathbf{A}}\) be the unital stellar algebra deduced from A by adjoining a unit element \(e\) . Denote by \(\pi\) the hermitian character \(x + \lambda e \mapsto \lambda\) of \(\widetilde{\mathbf{A}}\) into \(\mathbf{C}\) ; we have \(\operatorname{Ker}(\pi) = \mathbf{A}\) .

Let \(x \in \mathrm{A}\) be a normal element. It is normal in \(\widetilde{\mathrm{A}}\) and \(\mathrm{Sp}_{\widetilde{\mathrm{A}}}^{\prime}(x) = \mathrm{Sp}_{\mathrm{A}}^{\prime}(x)\) . Denote by \(\mathcal{C}'(\mathrm{Sp}_{\mathrm{A}}^{\prime}(x))\) the stellar algebra of continuous functions \(f\) on \(\mathrm{Sp}_{\mathrm{A}}^{\prime}(x)\) such that \(f(0) = 0\) .

Let \(f \in \mathcal{C}(\mathrm{Sp}_{\mathrm{A}}^{\prime}(x))\) . Since \(\pi(f(x)) = f(\pi(x))\) (prop. 8 of I, p.~112), we have \(f(x) \in \mathrm{A}\) if and only if \(f(0) = 0\) . The map \(f \mapsto f(x)\) defines a morphism of involutive algebras from the stellar algebra \(\mathcal{C}'(\mathrm{Sp}_{\mathrm{A}}^{\prime}(x))\) to A for which the image of the identity map \(z\) of \(\mathrm{Sp}_{\mathrm{A}}^{\prime}(x)\) is \(x\) . This morphism is isometric and its image is the stellar subalgebra of A generated by \(x\) .

PROPOSITION 12. --- The map \(f \mapsto f(x)\) is the unique morphism of involutive algebras from the stellar algebra \(\mathcal{C}'(\mathrm{Sp}_{\mathrm{A}}'(x))\) to A such that the identity map \(z\) of \(\mathrm{Sp}_{\mathrm{A}}'(x)\) has image \(x\) .

The elements \(z\) and \(\overline{z}\) of \(\mathcal{C}'(\mathrm{Sp}_{\mathrm{A}}'(x))\) generate a dense subalgebra of \(\mathcal{C}'(\mathrm{Sp}_{\mathrm{A}}'(x))\) (cf.~TG, X, p.~40, cor. 2). Since every morphism of involutive algebras from the stellar algebra \(\mathcal{C}'(\mathrm{Sp}_{\mathrm{A}}'(x))\) to the stellar algebra A is continuous (I, p.~104, prop. 2), the result follows.

The results of the preceding number concerning functional calculus extend to the general case. We shall simply state them and leave it to the readers to complete the proofs, mutatis mutandis.

PROPOSITION 13. --- One has the following properties:

\begin{enumerate}
\def\labelenumi{\alph{enumi})}
\item
  For every \(f \in \mathcal{C}'(\mathrm{Sp}_{\mathrm{A}}'(x))\) , one has \(\mathrm{Sp}_{\mathrm{A}}'(f(x)) = f(\mathrm{Sp}_{\mathrm{A}}'(x))\) ;
\item
  For every \(f \in \mathcal{C}'(\mathrm{Sp}_{\mathrm{A}}'(x))\) and for every \(g \in \mathcal{C}'(\mathrm{Sp}_{\mathrm{A}}'(f(x)))\) , one has \((g \circ f)(x) = g(f(x))\) ;
\item
  Let A' be a stellar algebra and let \(\pi\) be a morphism from A to A'; then \(\pi(x)\) is normal in A', one has \(\mathrm{Sp}_{\mathrm{A}'}^{\prime}(\pi(x)) \subset \mathrm{Sp}_{\mathrm{A}}^{\prime}(x)\) and \(\pi(f(x)) = f(\pi(x))\) for every \(f \in \mathcal{C}'(\mathrm{Sp}_{\mathrm{A}}^{\prime}(x))\);
\item
  If A is commutative, and if \(f \in \mathcal{C}'(\mathrm{Sp}_{\mathrm{A}}'(x))\) , then \(\mathcal{G}_{\mathrm{A}}'(f(x)) = f \circ \mathcal{G}_{\mathrm{A}}'(x)\) .
\end{enumerate}

Remark. --- Let A be a unital stellar algebra and let \(\widetilde{\mathrm{A}}\) be the unital stellar algebra deduced from A by adjoining a unit element \(e\) . For every \(x \in \mathrm{A}\) , one has \(\mathrm{Sp}_{\mathrm{A}}^{\prime}(x) = \mathrm{Sp}_{\mathrm{A}}(x) \cup \{0\}\) . Let \(x\) be a normal element of A. It is then a normal element of \(\widetilde{\mathrm{A}}\) and one therefore has two functional calculus maps in A, the first defined on \(\mathcal{C}(\mathrm{Sp}_{\mathrm{A}}(x))\) and the second on \(\mathcal{C}'(\mathrm{Sp}_{\mathrm{A}}^{\prime}(x))\) . Let \(f' \in \mathcal{C}'(\mathrm{Sp}_{\mathrm{A}}^{\prime}(x))\) ; if one denotes by \(f\) its restriction to \(\mathrm{Sp}_{\mathrm{A}}(x)\) , we then have \(f'(x) = f(x)\) .

\subsection{Positive elements in stellar algebras}

DEFINITION 6. --- Let A be a stellar algebra. An element x of A is said to be positive if it is Hermitian and if \(\mathrm{Sp}_{\mathrm{A}}^{\prime}(x) \subset \mathbf{R}_{+}\) . One denotes by \(\mathrm{A}_{+}\) the set of positive elements of A. It is a subset of \(\mathrm{A}_{h}\) .

We denote \(x \geqslant y\) if \(x - y \in \mathrm{A}_{+}\) .

If the stellar algebra A is unital, its identity element is positive.

If B is a stellar subalgebra of A, one has \(B_{+} = B \cap A_{+}\) (cor. of prop. 5 of I, p.~106).

If \(\pi\colon A\to B\) is a morphism of stellar algebras, then \(\pi(A_{+})\subset B_{+}\) .

Examples. --- 1) Let X be a locally compact space. In the stellar algebra \(\mathcal{C}_{0}(\mathrm{X})\) of continuous functions on X tending to 0 at infinity, resp. in the stellar algebra \(\mathcal{C}_{b}(\mathrm{X})\) of bounded continuous functions on X, a function \(f\) is a positive element if and only if it is real-valued and if \(f(x) \geqslant 0\) for all \(x \in \mathrm{X}\) (cf.~example 3 of I, p.~17).

\begin{enumerate}
\def\labelenumi{\arabic{enumi})}
\setcounter{enumi}{1}
\item
  Let A be a commutative stellar algebra. Let \(a\) be an element of A. Since \(\mathrm{Sp}_{\mathrm{A}}^{\prime}(x)\) is the union of \(\{0\}\) and of the image of the Gelfand transform \(\mathcal{G}(a)\) (prop. 6 of I, p.~37), the element \(a\) is positive if, and only if, the Gelfand transform \(\mathcal{G}(a)\) is a positive function.
\item
  Let E be a complex Hilbert space. An element x of the stellar algebra \(\mathcal{L}(\mathrm{E})\) (example 1 of I, p.~102) is positive if and only if it is a positive endomorphism of E in the sense of EVT, V, p.~45, def. 6 (prop. 8 of I, p.~138).
\end{enumerate}

Lemma 12. --- Let A be a unital stellar algebra and let \(x \in A\) be a Hermitian element.

\begin{enumerate}
\def\labelenumi{\alph{enumi})}
\item
  The element \(x\) is positive if and only if \(\left\| \|x\| \cdot 1 - x \right\| \leqslant \|x\|\) ;
\item
  If \(\| x\| \leqslant 1\) , then \(x\) is positive if and only if \(\| 1 - x\| \leqslant 1\) .
\item
  If x is positive, then 1 - x is positive if and only if \(\|x\| \leqslant 1\) ;
\item
  If x is positive and if \(y \in A_{+}\) commutes with x, then xy is positive.
\end{enumerate}

The element \(x\) is Hermitian, hence normal. By considering the stellar subalgebra generated by \(x\) , which is commutative, one reduces to the case where the algebra A is commutative, that is, to the case where \(A = \mathcal{C}_0(X)\) for a locally compact topological space X (th. 1 of I, p.~108). The first three assertions then follow immediately from example 1 above. Likewise, to prove assertion \(d\) ), one may consider the stellar subalgebra generated by \(x\) and \(y\) , which is commutative.

PROPOSITION 14. --- Let A be a stellar algebra. The set \(\mathrm{A}_{+}\) is a closed pointed salient convex cone in the real Banach space \(\mathrm{A}_{h}\) (EVT, II, p.~11).

Let \(\widetilde{\mathsf{A}}\) be the stellar algebra deduced from \(\mathsf{A}\) by adjoining an identity element. Since \(\mathsf{A}_{+} = \mathsf{A} \cap \widetilde{\mathsf{A}}_{+}\) , it suffices to prove the proposition for \(\widetilde{\mathsf{A}}\) . One may therefore suppose that \(\mathsf{A}\) has a unit element.

We have \(0 \in \mathrm{A}_{+}\) . For every \(\lambda \in \mathbf{R}_{+}^{*}\) and every \(x \in \mathrm{A}\) , one has \(\mathrm{Sp}_{\mathrm{A}}^{\prime}(\lambda x) = \lambda \mathrm{Sp}_{\mathrm{A}}^{\prime}(x)\) which implies that \(\mathrm{A}_{+}\) is a cone in the real Banach space \(\mathrm{A}_{h}\) .

To show that \(A_{+}\) is convex, it suffices to show that if x and y are positive, then \(x + y \geqslant 0\) (EVT, II, p.~11, prop. 10). By homothety, it suffices to prove that if \(x \geqslant 0\) and \(y \geqslant 0\) additionally satisfy \(\|x\| \leqslant 1\) , \(\|y\| \leqslant 1\) , then the element \(\frac{1}{2}(x + y)\) is positive. But we have

\[
\left\| 1 - \frac {1}{2} (x + y) \right\| \leqslant \frac {1}{2} \| 1 - x \| + \frac {1}{2} \| 1 - y \| \leqslant 1
\]

according to the assertion \(b)\) of Lemma 12, and this same assertion then shows that \(\frac{1}{2}(x + y)\) is positive.

Finally, the assertion \(a\) ) of Lemma 12 also implies that \(\mathrm{A}_{+}\) is closed.

Since \(A_{+}\) is a pointed cone in \(A_{h}\) , it is salient if and only if \(\mathrm{A}_{+}\cap(-\mathrm{A}_{+})\) is reduced to 0. But if \(x\in\mathrm{A}_{+}\cap(-\mathrm{A}_{+})\) , one has \(\mathrm{Sp}_{\mathrm{A}}^{\prime}(x)=\{0\}\) , hence \(\varrho(x)=0\) , and \(\|x\|=0\) . Since x is Hermitian (Lemma 7, (2) of I, p.~104), it follows that x=0.

Proposition 14 means that the relation \(x\geqslant y\) is an order relation on \(\mathrm{A}_h\) (EVT, II, p.~13, prop. 13).

PROPOSITION 15. — Let A be a stellar algebra. Let x be a normal element of A.

\begin{enumerate}
\def\labelenumi{\alph{enumi})}
\item
  Suppose that A is unital and let f be a continuous function from \(\mathrm{Sp}_{\mathrm{A}}(x)\) to C. For \(f(x)\) to be positive, it is necessary and sufficient that the image of f be contained in \(R_{+}\) ;
\item
  Let f be a continuous function from \(\mathrm{Sp}_{\mathrm{A}}^{\prime}(x)\) to C such that \(f(0)=0\) . For \(f(x)\) to be a positive element of A, it is necessary and sufficient that the image of f be contained in \(R_{+}\) .
\end{enumerate}

Assertion \(a\) ) follows from assertion \(a\) ) of Cor. 2 of I, p.~111, and assertion \(b\) ) follows from Prop. 13 of I, p.~114.

Let \(x\) be a Hermitian element of the stellar algebra A. Its spectrum is contained in \(\mathbf{R}\) (Prop. 4 of I, p.~106). Consider the continuous functions from \(\mathrm{Sp}_{\mathrm{A}}^{\prime}(x)\) into \(\mathbf{R}\) defined by

\[
f _ {1} \colon t \mapsto \sup (t, 0), \quad f _ {2} \colon t \mapsto \sup (- t, 0), \quad f _ {3} \colon t \mapsto | t |.
\]

We denote

\[
x ^ {+} = f _ {1} (x), \quad x ^ {-} = f _ {2} (x), \quad | x | = f _ {3} (x).\tag{3}
\]

Since the functions \(f_1, f_2, f_3\) are real-valued, positive, and vanish at 0, the elements \(x^+, x^-\) and \(|x|\) are positive elements of A (Prop. 15, a)) which belong to the stellar subalgebra of A generated by \(x\) .

We have \(f_{1}(t) - f_{2}(t) = t\) for all \(t \in \mathbf{R}\) , as well as the relations \(f_{1} + f_{2} = f_{3}\) and \(f_{1}f_{2} = 0\) . It follows that

\[
x = x ^ {+} - x ^ {-}, \quad | x | = x ^ {+} + x ^ {-}, \quad x ^ {+} x ^ {-} = x ^ {-} x ^ {+} = 0.\tag{4}
\]

Since the functional calculus map is isometric, we have

\[
\| | x | \| = \| x \|, \quad \| x ^ {+} \| \leqslant \| x \|, \quad \| x ^ {-} \| \leqslant \| x \|.
\]

Let \(x\) be a positive element of A. It is Hermitian, hence normal. Let \(\alpha \in \mathbf{R}_{+}^{*}\) , and let \(g\) be the restriction to \(\mathrm{Sp}_{\mathrm{A}}^{\prime}(x)\) of the function \(t \mapsto t^{\alpha}\) ; we denote \(x^{\alpha} = g(x)\) . This is a positive element of the stellar subalgebra of A generated by \(x\) . Let \(\alpha\) and \(\beta\) be in \(\mathbf{R}_{+}^{*}\) . Since the functional calculus map is an algebra homomorphism, and by Prop. 13 of I, p.~114, we have

\[
x ^ {\alpha} x ^ {\beta} = x ^ {\alpha + \beta} \quad (x ^ {\alpha}) ^ {\beta} = x ^ {\alpha \beta}.\tag{5}
\]

We shall also write \(\sqrt{x} = x^{1/2}\) .

PROPOSITION 16. — Let x be a positive element of A. Let \(\alpha \in \mathbf{R}_{+}^{*}\) . Then \(x^{1/\alpha}\) is the unique positive element y of A such that \(y^{\alpha} = x\) .

We saw above that \(x^{1/\alpha}\) satisfies the required properties. Conversely, let y be a positive element of A such that \(y^{\alpha} = x\) . By formula (5), we have \(y = (y^{\alpha})^{1/\alpha} = x^{1/\alpha}\) , which is what had to be proved.

Lemma 13. — Let A be a unital stellar algebra. Every element of A is a sum of unitary elements.

Let \(x\) be a Hermitian element of A. Suppose first that \(\| x \| \leqslant 2\) . By Lemma 12, c), we have \(1 - \frac{1}{4} x^2 \in \mathrm{A}_+\) . Let \(y = \frac{1}{2} x + i\sqrt{1 - \frac{1}{4} x^2}\) . We have \(y^* = \frac{1}{2} x - i\sqrt{1 - \frac{1}{4} x^2}\) , hence \(yy^* = 1\) and \(x = y + y^*\) is a sum of two unitary elements. In the general case, let \(k\) be an integer such that \(\| \frac{1}{k} x \| \leqslant 2\) ; the element \(x\) is then a sum of \(2k\) unitary elements. By Lemma 2 of I, p.~96, the lemma follows.

THEOREM 2. — Let A be a stellar algebra. An element x of A is positive if and only if there exists \(y \in A\) such that \(x = y^{*}y\) .

Suppose that \(x\) is positive. Let \(y = x^{1/2}\) ; this is a Hermitian element of A and we have \(y^*y = y^2 = x\) .

Conversely, let y be an element of A and set \(x = y^{*}y\) . This is a Hermitian element of A. Let us show that x is positive. For this, denote \(z = x^{-}\) and set w = yz. We then have

\[
w ^ {*} w = z ^ {*} y ^ {*} y z = z x z = z (x ^ {+} - z) z = - z ^ {3}.
\]

Since \(z \geqslant 0\) , one has \(z^3 \geqslant 0\) and we deduce that \(\mathrm{Sp}_{\mathrm{A}}'(w^{*}w) \subset \mathbf{R}_{-}\) . Let us also write \(w = w_1 + iw_2\) where \(w_1\) and \(w_2\) are Hermitian (lemma 2 of I, p.~96). The elements \(w_1^2\) and \(w_2^2\) are positive. We have \(ww^{*} + w^{*}w = 2w_{1}^{2} + 2w_{2}^{2}\) , and prop. 14 therefore shows that \(ww^{*} = 2w_{1}^{2} + 2w_{2}^{2} + (-w^{*}w)\) is positive. Since \(\mathrm{Sp}_{\mathrm{A}}'(ww^{*}) = \mathrm{Sp}_{\mathrm{A}}'(w^{*}w)\) (I, p.~5, prop. 1), which is contained in \(\mathbf{R}_{-}\) , we conclude that \(\mathrm{Sp}_{\mathrm{A}}'(w^{*}w) = \{0\}\) . Since \(w^{*}w\) is Hermitian, this implies (cor. 1 of I, p.~108) that \(\|w^{*}w\| = \varrho(w^{*}w) = 0\) , so that \(z^{3} = 0\) . Since z is Hermitian, we have z = 0. Thus, \(x = x^{+}\) is positive.

Remark. --- Let A be a stellar algebra and let \(x \in A\) . The element \(x^{*}x\) is positive; we therefore set \(|x| = (x^{*}x)^{1/2}\) . We have \(\|x\|^{2} = \|x^{*}x\| = |||x|^{2}\| = |||x||^{2}\) , hence \(\|x\| = |||x|\|\) .

When x is normal, we also have \(|x| = f(x)\) , where f is the map from C to C, vanishing at 0, given by \(f(z) = |z|\) . In particular, when x is Hermitian, \(|x|\) coincides with the element defined by formula (3).

Suppose further that A is unital and that x is invertible. Then \(|x|\) is also invertible, the element \(u = x|x|^{-1}\) is unitary and we have \(x = u|x|\) (polar decomposition; see also I, p.~139, no. 8 for the case of endomorphisms of Hilbert spaces).

Lemma 14. --- Let A be a stellar algebra, and let x and y be Hermitian elements of A.

\begin{enumerate}
\def\labelenumi{\alph{enumi})}
\item
  If \(x \leqslant y\) , then for every element \(w\) of A, we have \(w^{*}xw \leqslant w^{*}yw\) . In particular, if \(y \geqslant 0\) , one has \(w^{*}yw \geqslant 0\) ;
\item
  Suppose that A is unital. If \(0 \leqslant x \leqslant y\) and if \(x\) is invertible, then \(y\) is invertible and \(y^{-1} \leqslant x^{-1}\) ;
\item
  If \(0 \leqslant x \leqslant y\) then \(\|x\| \leqslant \|y\|\) .
\end{enumerate}

Let \(u = (y - x)^{1/2}\) . We have \(w^{*}yw - w^{*}xw = w^{*}u^{2}w = (uw)^{*}(uw)\) and the assertion a) follows from th. 2.

Let us prove assertion \(b\) ). Suppose first that \(x = 1\) . Let B be the unital stellar subalgebra generated by \(y\) . By the Gelfand isomorphism, \(y\) corresponds to a continuous function \(\geqslant 1\) on the compact space \(\mathsf{X}(\mathsf{B})\) . This function is therefore invertible, and its inverse is \(\leqslant 1\) . This implies that \(y\) is invertible and \(y^{-1} \leqslant 1 = x^{-1}\) . In the general case, we observe that \(0 \leqslant 1 \leqslant x^{-1/2}yx^{-1/2}\) according to \(a\) ), so the preceding case implies that \(z = x^{-1/2}yx^{-1/2}\) is invertible and \(z^{-1} \leqslant 1\) . Thus, \(y\) is invertible and \(y^{-1} \leqslant x^{-1}\) according to \(a\) ) again.

To prove the assertion \(c\) ), we may assume that A is unital (prop. 3 of I, p.~105). Suppose first that \(y\) is invertible. Let \(b = \sqrt{y}\) . According to \(a\) ), the conditions \(0 \leqslant x \leqslant y\) imply \(0 \leqslant b^{-1}xb^{-1} \leqslant b^{-1}yb^{-1} = 1\) , hence \(\|b^{-1}xb^{-1}\| \leqslant 1\) by Lemma 12, \(c\) ) of I, p.~116. We then have

\[
\| x \| = \| b (b ^ {- 1} x b ^ {- 1}) b \| \leqslant \| b \| \| b ^ {- 1} x b ^ {- 1} \| \| b \| \leqslant \| b \| ^ {2} = \| b ^ {2} \| = \| y \|.
\]

In the general case, for every real \(\varepsilon > 0\) , the element \(y + \varepsilon\) is invertible and \(0 \leqslant x \leqslant y + \varepsilon\) . From what precedes, we therefore have \(\|x\| \leqslant \|y + \varepsilon\|\) for every real number \(\varepsilon > 0\) , whence the result.

Note. --- In general, if x and y are positive elements of a stellar algebra A, the condition \(0 \leqslant x \leqslant y\) does not imply \(x^{2} \leqslant y^{2}\) (cf. exercise 15 of I, p. 182).

\subsection{Approximate units in stellar algebras}

DEFINITION 7. --- Let A be a normed algebra. An approximate unit of A is a filter base F on the unit ball of A such that, for every x in A, the filter bases xF and Fx on A converge to x, in other words:

\[
\lim _ {f, \mathfrak {F}} f x = \lim _ {f, \mathfrak {F}} x f = x.
\]

If A is a stellar algebra, an approximate unit F is said to be increasing if F is a filter base on \(A_{+}\) .

Let A be a stellar algebra. We denote by \(A_{+}^{\leqslant1}\) (resp. \(A_{+}^{<1}\) ) the set of positive elements of A of norm \(\leqslant 1\) (resp. of norm \(< 1\)); these are the Hermitian elements of A whose spectrum is contained in \([0,1]\) (respectively in \([0,1)\)).

PROPOSITION 17. --- Let A be a stellar algebra, let \(\mathfrak{F}\) be a filter base on \(\mathrm{A}_{+}^{\leqslant 1}\) . For \(\mathfrak{F}\) to be an increasing approximate unit of A, it is necessary and sufficient that one have

\[
\lim _ {f, \mathfrak {F}} f x = x
\]

for every positive element x of A.

The condition is obviously necessary; let us prove that it is sufficient. Let \(\widetilde{\mathrm{A}}\) be the stellar algebra deduced from A by adjoining a unit element. Let \(x\) be an element of A and let \(f \in \mathrm{A}_{+}^{\leqslant 1}\) . We have

\[
\begin{array}{c} \| f x - x \| ^ {2} = \| (f x - x) (f x - x) ^ {*} \| = \| (f - 1) x x ^ {*} (f - 1) \| \\ \leqslant \| (f - 1) x x ^ {*} \| \end{array}
\]

because \(\| f - 1\| \leqslant 1\) (Lemma 12, \(c\) ) of I, p.~116). We therefore have

\[
\limsup _ {f, \mathfrak {F}} \| f x - x \| ^ {2} \leqslant \limsup _ {f, \mathfrak {F}} \| f x x ^ {*} - x x ^ {*} \|.
\]

Since \(xx^{*}\) is positive (th. 2 of I, p.~118), the hypothesis implies that

\[
\limsup _ {f, \mathfrak {F}} \| f x - x \| ^ {2} \leqslant \| x x ^ {*} - x x ^ {*} \| = 0,
\]

so that

\[
\lim _ {f, \mathfrak {F}} f x = x.
\]

Since the involution of A is continuous and since \(\mathfrak{F}\) is a filter base on \(A_{h}\) , it follows that

\[
\lim _ {f, \mathfrak {F}} x f = \lim _ {f, \mathfrak {F}} (f ^ {*} x ^ {*}) ^ {*} = \lim _ {f, \mathfrak {F}} (f x ^ {*}) ^ {*} = (x ^ {*}) ^ {*} = x.
\]

The proposition follows from this.

PROPOSITION 18. --- Let A be a stellar algebra. The ordered set \(A_{+}^{<1}\) is directed to the right (E, III, p.~12, def. 7) and the filter of its sections (TG, I, p.~38, example 2) is an increasing approximate unit of A.

Let \(\widetilde{\mathsf{A}}\) be the stellar algebra deduced from A by adjoining a unit element denoted 1 (def. 4 of I, p.~106).

Let \(g\) be the function on \([0,1)\) into \(\mathbf{R}_+\) defined by \(g(t) = t(1 - t)^{-1}\) . This is a continuous increasing bijection, whose inverse bijection is given by \(t \mapsto 1 - (1 + t)^{-1}\) .

Let us prove that the ordered set \(A_{+}^{<1}\) is directed to the right. Let x and y be elements of \(A_{+}^{<1}\) . Since \(g(0)=0\) and since \(\mathrm{Sp}_{\mathrm{A}}^{\prime}(x)\) and \(\mathrm{Sp}_{\mathrm{A}}^{\prime}(y)\) are contained in \([0,1)\) , the elements \(g(x)\) and \(g(y)\) are defined; they are positive, so that \(g(x)+g(y)\geqslant0\) . We can therefore form the element \(z=g^{-1}(g(x)+g(y))\) of A. We have \(\mathrm{Sp}_{\mathrm{A}}^{\prime}(z)\subset[0,1)\) , and hence \(z\in A_{+}^{<1}\) .

We have \(0 \leqslant g(x) \leqslant g(z)\) , whence \(1 \leqslant 1 + g(x) \leqslant 1 + g(z)\) . Assertion \(b\) ) of Lemma 14 of I, p.~119 implies that \(1 + g(x)\) and \(1 + g(z)\) are invertible and that \((1 + g(z))^{-1} \leqslant (1 + g(x))^{-1}\) . Consequently, \(z = 1 - (1 + g(z))^{-1} \geqslant 1 - (1 + g(x))^{-1} = x\) . Similarly, we have \(z \geqslant y\) . Therefore, \(z\) majorizes \(x\) and \(y\) in \(\mathrm{A}_{+}^{<1}\) . The ordered set \(\mathrm{A}_{+}^{<1}\) is thus directed to the right. Let \(\mathfrak{F}\) denote the filter of its sections

Let \(x\) be a positive element of A. For every integer \(n \geqslant 1\) , let \(e_n = g^{-1}(nx)\) ; we have \(e_n \in \mathrm{A}_+^{<1}\) . Let \(h_n\) be the continuous function on \(\mathbf{R}_+\) defined for all \(t \in \mathbf{R}_+\) by \(h_n(t) = t^2(1 - g^{-1}(nt)) = t^2 / (1 + nt)\) . We have \(|h_n(t)| \leqslant t / n\) for all \(t \geqslant 0\) , and hence \(\| x(1 - e_n)x \| = \| h_n(x) \| \leqslant \| x \| / n\) . In particular, \(x(1 - e_n)x\) tends to 0 as \(n\) tends to infinity.

Let \(\varepsilon > 0\) be a real number. Let \(n\) be an integer such that \(\|x(1 - e_n)x\| < \varepsilon\) . For every \(f \in \mathrm{A}_+^{<1}\) such that \(f \geqslant e_n\) , we then have

\[
\begin{array}{r l} & {\| x - f x \| ^ {2} = \| (1 - f) x \| ^ {2} = \| ((1 - f) x) ^ {*} (1 - f) x \| = \| x ^ {*} (1 - f) ^ {2} x \|} \\ & {\qquad \qquad \qquad \qquad \qquad \qquad \qquad \qquad \qquad \qquad = \| x (1 - f) ^ {2} x \|.} \end{array}
\]

Additionally, since \(0 \leqslant f \leqslant 1\) , it follows that \((1-f)-(1-f)^{2} = (1-f)f \geqslant 0\) (Lemma 12, d) of I, p.~116), hence \((1-f)^{2} \leqslant 1-f\) . Since \(1-f \leqslant 1-e_{n}\) , it follows that \(0 \leqslant (1-f)^{2} \leqslant 1-e_{n}\) . By lemma 14, a) and c) of I, p.~119, it follows

\[
\left\| x - f x \right\| ^ {2} = \left\| x (1 - f) ^ {2} x \right\| \leqslant \left\| x (1 - e _ {n}) x \right\| <   \varepsilon .
\]

We therefore have \(\lim_{f,\mathfrak{F}}fx = x\) for all \(x \in A_{+}\) . The filter \(\mathfrak{F}\) is therefore an approximate unit of A by proposition 17.

\subsection{Quotient by a closed two-sided ideal}

Lemma 15. --- Let A be a stellar algebra and let I be a closed two-sided ideal of A. Then I is self-adjoint.

Let \(J = I \cap I^*\) . The set \(J\) is a self-adjoint two-sided ideal of \(A\) , which contains \(I^*I\) . In particular, \(J\) is a stellar algebra. Let \(\mathfrak{F}\) be an increasing approximate unit of \(J\) (prop. 18 of I, p.~121). For all \(x \in I\) and every \(f \in J_{+}^{\leqslant 1}\) , one has

\[
\begin{array}{c} \| x f - x \| ^ {2} = \| (x f - x) ^ {*} (x f - x) \| = \| f (x ^ {*} x f - x ^ {*} x) - (x ^ {*} x f - x ^ {*} x) \| \\ \leqslant 2 \| x ^ {*} x f - x ^ {*} x \|. \end{array}
\]

Since \(x^{*}x\in \mathrm{J}\) , we therefore have

\[
\lim _ {f, \mathfrak {F}} \| x f - x \| ^ {2} = 0.
\]

Since \(xf \in \mathrm{J}\) for all \(f \in \mathrm{J}_{+}^{\leqslant 1}\) and since \(\mathrm{J}\) is closed, this implies that \(x \in \mathrm{J}\) . Therefore \(\mathrm{I} = \mathrm{J}\) , and the ideal \(\mathrm{I}\) is self-adjoint.

PROPOSITION 19. --- Let A be a stellar algebra and let I be a closed two-sided ideal of A. Then the quotient involutive Banach algebra A/I is a stellar algebra.

The ideal I is self-adjoint (Lemma 15). We consider the stellar algebra \(\widetilde{\mathsf{A}}\) deduced from A by adjoining an identity element (def. 4 of I, p.~106). In this algebra, the set I is a closed self-adjoint two-sided ideal, and A/I is identified with a closed involutive subalgebra of \(\widetilde{\mathrm{A}}/\mathrm{I}\) . We may therefore assume that A is unital.

The Banach algebra A/I is involutive. Let \(\pi\colon\mathrm{A}\to\mathrm{A}/\mathrm{I}\) be the canonical projection. The self-adjoint two-sided ideal I is a stellar subalgebra of A. Let \(\mathfrak{F}\) be an increasing approximate unit of I (prop. 18 of I, p.~121). Let us first show that for every \(x\in\mathrm{A}\) , one has

\[
\| \pi (x) \| _ {\mathrm{A/I}} = \lim _ {f, \mathfrak {F}} \| x - x f \|.\tag{6}
\]

On the one hand, since \(xf \in \mathrm{I}\) for all \(f \in \mathrm{I}_{+}^{\leqslant 1}\) , one has

\[
\| \pi (x) \| _ {\mathrm{A} / \mathrm{I}} = \inf _ {a \in \mathrm{I}} \| x - a \| \leqslant \operatorname * {l i m i n f} _ {f, \mathfrak {F}} \| x - x f \|.
\]

On the other hand, for every \(a \in \mathrm{I}\) , one has

\[
\begin{array}{c} \| x - x f \| \leqslant \| (x - a) - (x - a) f \| + \| a - a f \| \\ = \| (x - a) (1 - f) \| + \| a - a f \| \end{array}
\]

and therefore, since \(\|1 - f\| \leqslant 1\) (Lemma 12 of I, p.~116) and \(a \in I\), from which we deduce

\[
\operatorname * {l i m s u p} _ {f, \mathfrak {F}} \| x - x f \| \leqslant \| x - a \|.
\]

Thus, it follows \(\limsup_{f,\mathfrak{F}} \|x - xf\| \leqslant \|\pi(x)\|_{\mathrm{A/I}}\) since \(a\) is arbitrary in I. Formula (6) is therefore proved.

Now let \(x\) be an element of A. By formula (6), we have

\[
\begin{array}{c} \| \pi (x) \| _ {\mathrm{A/I}} ^ {2} = \lim _ {f, \mathfrak {F}} \| x - x f \| ^ {2} = \lim _ {f, \mathfrak {F}} \| x (1 - f) \| ^ {2} \\ = \lim _ {f, \mathfrak {F}} \| (1 - f) x ^ {*} x (1 - f) \| \leqslant \lim _ {f, \mathfrak {F}} \| x ^ {*} x (1 - f) \| = \| \pi (x ^ {*} x) \| _ {\mathrm{A/I}}. \end{array}
\]

Lemma 6 of I, p.~103 then implies that the involutive Banach algebra A/I is a stellar algebra.

\subsection{Enveloping stellar algebra of an involutive Banach algebra}

Lemma 16. --- Let A be an involutive algebra and let p be a semi-norm on A. The following conditions are equivalent:

\begin{enumerate}
\def\labelenumi{(\roman{enumi})}
\item
  We have \(p(xy) \leqslant p(x)p(y)\) , \(p(x^{*}) = p(x)\) and \(p(x)^{2} = p(x^{*}x)\) for all \(x, y \in \mathrm{A}\) ;
\item
  The set R of elements x of A such that \(p(x) = 0\) is a self-adjoint two-sided ideal of A, and the seminorm on A/R induced by p makes A/R an involutive normed algebra whose completion is a C*-algebra;
\item
  There exists a stellar algebra B and a morphism of involutive algebras \(\varphi\) from A to B such that \(p(x) = \|\varphi(x)\|\) for all \(x\in\mathbf{A}\) .
\end{enumerate}

The implications (i) \(\Longrightarrow\) (ii) \(\Longrightarrow\) (iii) \(\Longrightarrow\) (i) are all elementary.

A semi-norm satisfying the conditions of Lemma 16 will be called a stellar semi-norm on the involutive algebra A.

Let A be a normed involutive algebra and let S be the set of stellar seminorms on A. We have \(p(x) \leqslant \|x\|\) for all \(x \in A\) and every \(p \in S\) (I, p.~104, prop. 2). The map \(x \mapsto \|x\|_* = \sup_{p \in S} p(x)\) is a stellar seminorm on A. It is the largest stellar seminorm on A.

Let R be the set of \(x \in A\) such that \(\|x\|_{*} = 0\) . It is a closed two-sided ideal of A. We denote by Stell(A) the completed stellar algebra of A/R for the norm induced by \(x \mapsto \|x\|_{*}\) (Lemma 16, (ii)). The canonical map from A into Stell(A) is continuous, its image is dense in Stell(A), and its kernel is equal to R.

DEFINITION 8. --- The stellar algebra Stell(A) is called the enveloping stellar algebra of the normed involutive algebra A.

If A is commutative, then Stell(A) is commutative; if A is unital, then Stell(A) is unital.

PROPOSITION 20. --- Let A be a normed involutive algebra and let j be the canonical morphism from A into Stell(A). For every stellar algebra B and for every morphism \(\varphi\) of involutive algebras from A to B, there exists a unique morphism \(\varphi'\) of stellar algebras from Stell(A) into B such that \(\varphi = \varphi' \circ j\) .

We denote by \(x \mapsto \|x\|_*\) the norm on Stell(A). Let R be the kernel of \(j\) . The map \(x \mapsto \| \varphi(x) \|\) is a stellar seminorm on A. We therefore have \(\| \varphi(x) \| \leqslant \| x \|_*\) for all \(x \in A\) . The morphism \(\varphi\) therefore defines, by passing to the quotient, a continuous morphism from A/R to B, which extends by continuity to a morphism \(\varphi'\) from Stell(A) to B satisfying \(\varphi = \varphi' \circ j\) . The uniqueness of \(\varphi'\) results from the fact that the image of \(j\) is dense in Stell(A).

COROLLARY. --- Let A be a commutative involutive Banach algebra and let j be the canonical morphism from A into Stell(A). The map \(\mathsf{X}(j)\) is a homeomorphism from \(\mathsf{X}(\mathrm{Stell}(\mathrm{A}))\) onto the subspace \(\mathsf{X}(\mathrm{A})_{h}\) of \(\mathsf{X}(\mathrm{A})\) consisting of the Hermitian characters of A.

The Hermitian characters of A are the morphisms of involutive algebras from A to the stellar algebra C. Proposition 20 therefore implies that \(\mathsf{X}(j)\) is a bijection from \(\mathsf{X}(\mathrm{Stell}(\mathsf{A}))\) onto \(\mathsf{X}(\mathsf{A})_h\) . Since \(\mathsf{X}(j)\) is a homeomorphism onto its image (cf.~I, p.~10), the corollary follows.

We identify \(\mathsf{X}(\mathrm{Stell}(\mathbf{A}))\) with \(\mathsf{X}(\mathbf{A})_h\) by the map \(\mathsf{X}(j)\) . For every \(x\in \mathbf{A}\) , the map \(\mathcal{G}_{\mathrm{Stell(A)}}(j(x))\) is none other than the restriction to \(\mathsf{X}(\mathbf{A})_h\) of \(\mathcal{G}_{\mathbf{A}}(x)\) .

PROPOSITION 21. --- Let A be an involutive Banach algebra and let j be the canonical morphism from A into Stell(A). The radical of A is contained in the kernel of j.

Let \(x\) be an element of the radical of A. Then \(x^{*}x\) belongs to the radical of A, and therefore \(\mathrm{Sp}_{\mathrm{A}}^{\prime}(x^{*}x) = \{0\}\) (I, p.~5, remark 3). Since \(\mathrm{Sp}_{\mathrm{Stell(A)}}^{\prime}(j(x^{*}x)) \subset \mathrm{Sp}_{\mathrm{A}}^{\prime}(x^{*}x)\) , we therefore have \(\mathrm{Sp}_{\mathrm{Stell(A)}}^{\prime}(j(x)^{*}j(x)) = \{0\}\) , whence \(\| j(x)\| ^2 = \| j(x)^{*}j(x)\| = \varrho (j(x)^{*}j(x)) = 0\) (formula (2) of I, p.~104), and therefore \(j(x) = 0\) .
